
\medterm{Fab Immunoglobulins} The target-binding part of an antibody. A laboratory-made Fab fragment can recognize and block a harmful substance or help remove it from the bloodstream in selected treatments.

\medterm{Fabry Disease} An inherited disorder in which low activity of an enzyme causes a fatty substance to accumulate in cells. It may cause burning pain in the hands and feet, skin spots, reduced sweating, kidney disease, heart disease, or stroke.

\medterm{Fabry'S Disease} An inherited enzyme disorder that causes a fatty substance to accumulate in cells. It may cause burning limb pain, skin findings, digestive symptoms, kidney or heart disease, and blood-vessel complications.

\medterm{Face} The front part of the head, including the eyes, nose, mouth, cheeks, and jaw. Its muscles produce facial expression, its nerves provide movement and sensation, and its structures support breathing, eating, speech, and vision.

\medterm{Face Pain} Pain in the face that may arise from teeth, sinuses, nerves, eyes, jaw joints, skin, blood vessels, or other tissues. Location, duration, triggers, swelling, fever, and neurologic signs help identify the cause.

\medterm{Face Powder Poisoning} Illness after swallowing or inhaling face powder or its ingredients. Coughing, choking, vomiting, breathing difficulty, drowsiness, or persistent symptoms require poison-control or emergency advice; do not induce vomiting.

\medterm{Face Presentation} A birth position in which the baby’s neck is extended so the face, rather than the top of the head, leads through the birth canal. The baby’s chin direction helps determine whether vaginal birth is safe.

\medterm{Facelift} Cosmetic surgery that removes or repositions selected facial and neck skin and deeper tissues to reduce sagging. It does not stop aging and may cause bleeding, infection, nerve injury, scarring, or temporary swelling and numbness.

\synonyms
Rhytidectomy

\medterm{Facet Joint} A small joint at the back of the spine connecting neighboring vertebrae and guiding movement. Arthritis, injury, or inflammation can cause localized neck or back pain and stiffness.

\medterm{Facet Joints} Small paired joints at the back of the spine that guide bending and rotation and help stabilize vertebrae. Wear, inflammation, or injury may cause localized neck or back pain.

\medterm{Facet Rhizotomy} A procedure that uses heat or another method to interrupt small nerves carrying pain signals from joints at the back of the spine. It may reduce selected chronic neck or back pain, but relief may be temporary.
\medterm{Facial Asymmetry} A difference in the size, position, movement, or shape of the two sides of the face. It may be present from birth or result from injury, nerve weakness, muscle disease, swelling, or a growth.

\medterm{Facial Cleft} A birth defect caused by incomplete joining of facial tissues during development. It may affect the lip, nose, mouth, or other facial structures and can interfere with feeding, speech, breathing, hearing, or appearance.

\medterm{Facial Edema} Swelling of the face caused by excess fluid in tissues. Allergy, infection, injury, kidney or heart disease, and blocked veins may cause it; sudden lip, tongue, or throat swelling with breathing difficulty is an emergency.

\medterm{Facial Expression} The visible movement or position of facial muscles that communicates emotion, pain, intention, or other information. Reduced or unequal expression can indicate facial-nerve or muscle dysfunction.

\medterm{Facial Hair, Excess In Women} Coarse dark hair on the face or other androgen-sensitive areas in a woman, called hirsutism. It may result from polycystic ovary syndrome, medicines, or rarely an adrenal or ovarian disorder; sudden rapid growth needs assessment.

\synonyms
Hirsutism

\medterm{Facial Muscle} A muscle that moves the skin of the face for expression and helps with speaking, eating, breathing through the mouth, or closing the eyes. Weakness may result from nerve, muscle, or brain disease.

\medterm{Facial Neoplasm} A benign or malignant growth arising in facial skin, soft tissue, bone, salivary glands, nerves, or other structures. Symptoms and treatment depend on its tissue of origin, size, location, and spread.

\medterm{Facial Nerve} The seventh cranial nerve, which controls facial expression and carries taste from part of the tongue. It also contributes to tear and saliva production and helps protect the ear from loud sounds.

\medterm{Facial Nerve Palsy} Weakness or paralysis of facial muscles caused by dysfunction of the seventh cranial nerve. It may affect eye closure, smiling, taste, tears, or saliva; sudden weakness with arm or speech symptoms may indicate a stroke.

\medterm{Facial Nerve Palsy Due To Birth Trauma} One-sided facial weakness in a newborn caused by pressure on the facial nerve during birth. It often improves as the nerve recovers, but eye protection, feeding, and other causes of weakness must be assessed.

\medterm{Facial Nerve Problems} Disorders affecting the nerve that controls facial movement and some taste, tear, saliva, and sound-protection functions. Causes include Bell palsy, stroke, infection, injury, or a tumor; sudden weakness needs emergency assessment.

\synonyms
Facial nerve disorder; facial nerve dysfunction

\medterm{Facial Nucleus} A group of nerve cells in the pons, part of the brainstem, whose fibers form the facial nerve and control facial-expression muscles. Injury can cause weakness of the face and related functions.

\medterm{Facial Pain} Pain in the face that may arise from dental or sinus disease, jaw-joint problems, nerve disorders, infection, injury, or other causes. The pattern and associated symptoms guide evaluation.

\medterm{Facial Palsy} Weakness or paralysis of facial muscles, usually on one side, caused by facial-nerve dysfunction. Bell palsy is one cause, but stroke, infection, injury, Lyme disease, and tumors are also possible; sudden onset with other neurologic signs is urgent.

\synonyms
Bell's palsy; facial paralysis

\medterm{Facial Paralysis} Partial or complete loss of facial-muscle movement caused by disease or injury of the facial nerve or its brain pathways. Sudden onset, especially with limb weakness or speech trouble, requires urgent stroke assessment.

\medterm{Facial Paresis} Weakness of muscles supplied by the facial nerve, causing reduced facial movement and sometimes poor eye closure, altered taste, or abnormal tearing. The cause may be nerve inflammation, injury, infection, stroke, or a tumor.

\medterm{Facial Spasm} Involuntary repeated tightening of facial muscles. It may affect one side around the eye or mouth and can result from nerve irritation, medicines, stress, or another neurologic disorder; persistent spasms need assessment.

\medterm{Facial Swelling} Enlargement of part or all of the face caused by infection, allergy, injury, dental or salivary disease, sinus disease, fluid retention, or a growth. Rapid swelling of the mouth or throat with breathing difficulty is an emergency.

\medterm{Facial Tics} Sudden, repeated, involuntary facial movements or sounds, such as blinking, grimacing, or nose wrinkling. They may increase with stress and occur in tic disorders; persistent or disruptive tics can be treated.

\medterm{Facial Trauma} Injury affecting facial skin, soft tissue, teeth, nerves, blood vessels, eyes, nose, or bones. Assessment must consider airway, bleeding, vision, brain injury, and structural damage.

\medterm{Facial Vein} A vein that drains blood from the superficial face toward the internal jugular vein. Connections with deeper veins mean facial infections can rarely spread toward the eye socket or brain.

\medterm{Facies} A characteristic facial appearance associated with a particular disease, inherited condition, or physiologic state. It can provide clues but is not diagnostic by itself and must be interpreted with other findings.
\medterm{Facilitated Diffusion} Movement of a substance across a cell membrane through a channel or carrier protein, from higher to lower concentration. The process does not directly use cellular energy and helps move substances that cannot cross the membrane unaided.

\medterm{Facioscapulohumeral Dystrophy} An inherited muscle disease causing gradually progressive weakness of the face, shoulder blades, upper arms, and sometimes legs or trunk. Onset and severity vary, and breathing, hearing, and eye complications may occur.

\medterm{Facioscapulohumeral Muscular Dystrophy} An inherited disorder causing progressive, often uneven weakness of facial, shoulder-blade, and upper-arm muscles. Symptoms may begin in adolescence or adulthood and can later involve the trunk, hips, or legs.

\medterm{Factitious Disorder} A mental-health disorder in which a person deliberately creates, exaggerates, or falsely reports illness to be treated as sick, without a clear outside reward. Care should be nonjudgmental and avoid unnecessary tests or procedures.

\synonyms
Munchausen Syndrome

\medterm{Factitious Disorders} Factitious disorder involves deliberately producing or falsifying illness behavior to assume the sick role, without an obvious external reward. It is a mental-health disorder requiring nonjudgmental assessment, protection from unnecessary procedures, and coordinated psychiatric and medical care.

\medterm{Factitious Hyperthyroidism} A state of excessive thyroid hormone in the body caused by ingestion of thyroid hormone or another external source rather than increased thyroid-gland production.

\medterm{Factor H Deficiency} A rare inherited or acquired problem in which a protein that keeps part of the immune system under control is missing or faulty. It can cause repeated infections or inflammation that damages small blood vessels, especially in the kidneys.

\medterm{Factor} A blood protein needed to form a clot. Different clotting factors work together to stop bleeding; too little or abnormal activity can cause a bleeding disorder, while excessive clotting can contribute to thrombosis.

\medterm{Factor I Deficiency} A rare condition in which a protein that helps control immune-system activity is missing or does not work properly. It may cause repeated bacterial infections or inflammation and injury to the kidneys.

\medterm{Factor II (Prothrombin) Assay} A blood test measuring factor II activity, which helps assess the clotting pathway and investigate abnormal bleeding, prolonged prothrombin time, or suspected prothrombin deficiency or inhibition.

\medterm{Factor IX} A vitamin-K-dependent blood-clotting protein that works with factor VIII to form clots. Inherited deficiency causes hemophilia B, leading to prolonged or spontaneous bleeding, especially into joints and muscles.
\medterm{Factor IX Assay} A blood test measuring factor IX activity, used to diagnose or monitor factor IX deficiency, the cause of hemophilia B, and to guide replacement treatment.

\medterm{Factor V} A blood-clotting protein that helps generate thrombin and fibrin, the mesh that strengthens a clot. Too little can cause bleeding; an inherited change affecting its regulation can increase venous-clot risk.

\medterm{Factor V Assay} A blood test measuring how well factor V supports clotting. It helps investigate abnormal bleeding or clotting-test results and diagnose inherited deficiency or an acquired reduction from liver disease or other illness.

\medterm{Factor V Deficiency} An inherited or acquired bleeding disorder caused by too little active factor V. It may cause easy bruising, nosebleeds, heavy menstrual bleeding, or excessive bleeding after injury, dental work, surgery, or childbirth.

\medterm{Factor V Leiden} An inherited change in factor V that makes it resist a natural clot-limiting protein. It increases the risk of clots in veins, but the risk varies with family history, pregnancy, hormones, surgery, and other factors.

\medterm{Factor V Leiden Mutation} A factor V gene variant that reduces the normal response to activated protein C, a clot-limiting substance. It causes inherited tendency toward venous thrombosis, although many carriers never develop a clot.

\medterm{Factor VII} A liver-made blood-clotting protein that starts clot formation when tissue is injured. Inherited or acquired deficiency can cause easy bruising, nosebleeds, or serious bleeding after injury or procedures.

\medterm{Factor VII Assay} A blood test measuring factor VII activity. It helps investigate a prolonged prothrombin time and identify inherited deficiency, liver disease, vitamin-K deficiency, or medicine-related reduction.

\medterm{Factor VII Deficiency} A bleeding disorder caused by too little or poorly functioning factor VII. Bleeding severity varies and may include bruising, nosebleeds, heavy periods, bleeding after surgery, or severe hemorrhage; prothrombin time is usually prolonged.

\medterm{Factor VIII} A blood-clotting protein that works with factor IX to form a stable clot. Inherited deficiency causes hemophilia A, with prolonged or spontaneous bleeding, often into joints, muscles, or other tissues.

\medterm{Factor VIII Assay} A blood test measuring factor VIII activity. It helps diagnose or monitor hemophilia A and investigate acquired reductions caused by inhibitors, liver disease, or other conditions.

\medterm{Factor X Assay} A blood test measuring factor X activity, a key clotting step shared by two clotting pathways. It helps investigate abnormal bleeding and diagnose or monitor inherited or acquired factor X deficiency.

\medterm{Factor X Deficiency} A rare inherited or acquired bleeding disorder caused by too little active factor X. Bleeding severity varies from bruising to serious hemorrhage, and both major clotting-time tests may be prolonged.

\medterm{Factor XI Deficiency} An inherited or acquired bleeding disorder caused by reduced factor XI activity. Bleeding is variable and is often noticed after surgery, dental procedures, childbirth, or trauma rather than occurring spontaneously.

\medterm{Factor Xii} A blood-clotting protein that starts one pathway measured by laboratory clotting tests. Factor XII deficiency can greatly prolong the activated partial-thromboplastin time but usually does not cause abnormal bleeding.

\medterm{Factor XII (Hageman Factor) Deficiency} A deficiency of factor XII that can greatly prolong one laboratory clotting test but usually does not cause bleeding. It is often discovered unexpectedly during testing for another reason.

\medterm{Factor XII Assay} A blood test measuring factor XII activity. It may help explain an isolated prolonged activated partial-thromboplastin time, especially when the person has no unusual bleeding.

\medterm{Factor Xiii} The clot-stabilizing protein that strengthens the fibrin mesh after a clot forms. Factor XIII deficiency can cause delayed bleeding, poor wound healing, or repeated pregnancy loss despite initially normal routine clotting tests.





\medterm{Facultative Anaerobe} A microorganism that can grow with or without oxygen. It uses oxygen when available but can switch to other energy-producing processes when oxygen is absent; many intestinal bacteria have this ability.
\medterm{Faecal Fistula} An abnormal connection between the intestine and another organ, the skin, or the body surface that allows stool or intestinal fluid to leak. It can cause infection, dehydration, skin injury, and poor nutrition.

\medterm{Faecalibacterium} A group of oxygen-sensitive bacteria commonly found in the intestine. They can produce butyrate, a substance that supports the colon lining; their presence is studied in relation to gut health and inflammation.

\medterm{Faeces} The solid waste passed from the intestine through the anus, also called stool. Its color, consistency, frequency, and contents can provide clues about digestion, bleeding, infection, inflammation, or other illness.

\medterm{Fahr Disease} A rare disorder involving abnormal calcium deposits in parts of the brain. It may cause movement problems, seizures, cognitive changes, or psychiatric symptoms; similar deposits can also result from other metabolic or genetic conditions.

\medterm{Failed Induction Of Labor} Failure of labor to progress after medicines or procedures intended to start it. The decision to continue induction or perform cesarean birth depends on cervical change, contractions, time, the mother’s condition, and fetal well-being.

\medterm{Failed Intubation} Inability to place a breathing tube safely in the windpipe. It is an airway emergency requiring prompt help, continued oxygen delivery, an alternative airway method, and a plan to avoid repeated traumatic attempts.

\medterm{Failure To Thrive} Poor weight gain or growth in a child. Causes may include feeding difficulty, inadequate intake, poor absorption, chronic disease, developmental problems, or social stress; assessment considers the child’s growth pattern, diet, development, and health.

\medterm{Fainting} A brief loss of consciousness and muscle tone caused by temporarily reduced blood flow to the brain, followed by spontaneous recovery. Reflex episodes are common, but fainting during exercise, with chest pain, or without warning needs urgent assessment.

\synonyms
syncope.

\medterm{Fainting (Syncope)} A temporary loss of consciousness and posture from reduced brain blood flow, followed by recovery. Causes include reflex reactions, standing-related blood-pressure drop, medicines, dehydration, and heart disease; warning features require urgent evaluation.

\medterm{Falciform Ligament} A fold of tissue attaching the liver to the front abdominal wall and diaphragm. It marks the division between the liver’s right and left lobes and contains the remnant of a fetal umbilical vein.

\medterm{Falciparum Malaria} Malaria caused by Plasmodium falciparum parasites transmitted by infected mosquitoes. It can worsen rapidly and cause brain involvement, severe anemia, kidney injury, breathing failure, shock, or death, so urgent treatment is required.

\medterm{Fall} An unplanned movement to the ground or another lower level. Falls may result from poor balance, weakness, vision problems, medicines, dizziness, low blood pressure, or environmental hazards and can cause serious injury.

\medterm{Fall Prevention} Steps to reduce falling, including strength and balance training, medication review, vision and hearing checks, safe footwear, adequate lighting, removal of hazards, and support with walking when needed.

\medterm{Fall Risk} The likelihood of falling based on factors such as previous falls, balance, walking ability, strength, vision, medicines, dizziness, blood-pressure changes, cognition, and environmental hazards.

\medterm{Fall Screening} A brief assessment of previous falls, balance, walking, strength, medicines, vision, dizziness, cognition, and surroundings to identify fall risk and guide prevention or further evaluation.

\medterm{Fallen Arches} Flattening of the foot’s inner arch, also called flat feet. It may be flexible or rigid and may cause no symptoms or lead to foot, ankle, knee, or back discomfort; treatment depends on symptoms and cause.

\synonyms
Pes planus; flat feet

\medterm{Fallopian Tube} One of two tubes between the ovaries and uterus. An egg usually travels through a tube, where fertilization may occur, before reaching the uterus; blockage or scarring can affect fertility.

\medterm{Fallopian Tube Cancer} A rare cancer beginning in a fallopian tube. It may cause pelvic pain, a mass, abnormal or watery discharge, or no early symptoms; diagnosis, staging, and treatment resemble those for ovarian cancer.

\medterm{Fallopian Tube Torsion} Twisting of a fallopian tube, sometimes around a cyst or fluid-filled tube, that can cut off blood flow. Sudden severe pelvic pain, nausea, or vomiting requires urgent surgical assessment.

\medterm{Fallopian Tubes} The two tubes connecting the ovaries and uterus. They transport eggs and are common sites of fertilization; blockage, inflammation, scarring, pregnancy, or tumors can affect them.

\medterm{Fallopian-Tube Endometriosis} Endometriosis involving or surrounding a fallopian tube. It may cause pelvic pain, adhesions, or infertility, or be found unexpectedly; treatment depends on symptoms, disease extent, and fertility goals.

\medterm{Fallot'S Tetralogy} A birth defect combining four heart abnormalities that reduce blood flow to the lungs and may cause bluish skin, breathing difficulty, poor growth, or fainting. Specialist care and usually surgery are required.

\medterm{Fallot, Tetralogy Of} A congenital heart defect consisting of four abnormalities: ventricular septal defect, right ventricular outflow tract obstruction, overriding aorta, and right ventricular hypertrophy. Causes cyanosis due to right-to-left shunting. Treatment includes surgical repair.

\synonyms
Tetralogy of Fallot

\medterm{Falls} Unplanned events in which a person comes to rest on the ground or another lower surface. Falls may cause fractures or head injury and often involve balance, weakness, vision, medicines, dizziness, or environmental hazards.

\medterm{False Labor} Contractions that may be uncomfortable but do not progressively open and thin the cervix. They are often irregular and may lessen with rest, hydration, or a change in activity, unlike established labor.

\medterm{False Negative} A test result that says a condition is absent when it is actually present. It can occur because testing was too early, the sample was inadequate, the condition level was low, or the test was not sensitive enough.

\medterm{False Passages} Abnormal channels made through tissue during instrumentation, surgery, or injury instead of following the normal body pathway. They can cause bleeding, perforation, infection, or failure to place a tube correctly.

\medterm{False Positive} A test result that indicates a condition is present when it is not. It may result from cross-reaction, contamination, technical error, or testing in a group where the condition is uncommon.
\medterm{False Ribs} The lower ribs that do not attach directly to the breastbone. Their cartilage joins the cartilage above or ends freely, allowing flexibility but leaving them vulnerable to injury in some areas.

\medterm{False-Negative Result} A negative result when the tested condition is actually present. Early testing, an inadequate sample, prior treatment, low levels of the target, technical error, or limited test sensitivity can cause it.

\medterm{False-Positive Result} A positive result when the tested condition is actually absent. Cross-reaction, contamination, technical error, limited specificity, or testing in a low-risk population can produce it.

\medterm{Faltering Weight} Slower-than-expected weight gain or a downward change in a child’s growth pattern. Feeding difficulty, inadequate intake, poor absorption, chronic disease, developmental problems, or social stress may contribute and require assessment.

\medterm{Famciclovir} An antiviral medicine that changes into penciclovir in the body and limits multiplication of herpes viruses. It can treat selected shingles, cold-sore, or genital-herpes episodes but does not remove the virus permanently.

\medterm{Familial} Occurring more often within a family because of shared genes, environment, behaviors, or a combination. Familial does not necessarily mean that one specific gene causes the condition.

\medterm{Familial Adenomatous Polyposis} An inherited condition causing numerous precancerous polyps in the colon and rectum. Without preventive treatment, colorectal cancer is very likely, so lifelong specialist screening and risk-reducing surgery are often needed.

\medterm{Familial Combined Hyperlipidemia} An inherited tendency to have high LDL cholesterol, high triglycerides, or both, often affecting several relatives. It increases early cardiovascular risk and is managed with lifestyle measures and lipid-lowering treatment when indicated.

\medterm{Familial Dysautonomia} An inherited nerve disorder affecting sensation and automatic functions such as blood pressure, heart rate, sweating, tears, digestion, and swallowing. Symptoms may include dizziness, temperature problems, repeated lung infections, and poor growth.

\medterm{Familial Dysbetalipoproteinemia} An inherited disorder in which the body clears certain cholesterol- and triglyceride-rich particles poorly. It can cause high blood fats, yellowish skin deposits, and premature artery disease, especially with other metabolic risk factors.

\medterm{Familial Hypercholesterolaemia} An inherited disorder causing very high LDL cholesterol from childhood. Without effective treatment, cholesterol can build up in artery walls and cause early heart attack, stroke, or disease of other arteries.

\medterm{Familial Hypercholesterolemia} An inherited disorder causing persistently high LDL cholesterol because the body removes it from blood poorly. Family history, cholesterol levels, physical findings, and sometimes genetic testing support the diagnosis.

\medterm{Familial Hypertriglyceridemia} An inherited tendency toward high blood triglycerides, often worsened by diabetes, obesity, alcohol, or certain medicines. Very high levels can cause abdominal pain, eruptive skin bumps, and acute pancreatitis.

\medterm{Familial Hypocalciuric Hypercalcemia} An inherited condition causing lifelong, usually mild high blood calcium with low urinary calcium. People often have no symptoms, and recognizing it prevents unnecessary parathyroid surgery; specialist testing may confirm the diagnosis.

\medterm{Familial Lipoprotein Lipase Deficiency} A rare inherited disorder in which triglycerides cannot be broken down normally. It causes extremely high triglycerides, recurrent pancreatitis, abdominal pain, and yellow-red skin bumps and requires specialist dietary and medical care.

\medterm{Familial Mediterranean Fever} An inherited inflammatory disorder causing repeated short attacks of fever with abdominal, chest, or joint pain. Without treatment, ongoing inflammation can damage the kidneys; preventive medicine can reduce attacks and complications.

\medterm{Familial Multiple Lipomatosis} An inherited tendency to develop many soft, slow-growing, usually painless fatty lumps under the skin, often on the trunk and limbs. They are generally benign, but painful, rapidly growing, or unusual lumps need assessment.

\medterm{Familial Multiple Trichoepithelioma} An inherited condition causing many benign tumors from hair follicles, usually small skin-colored facial bumps. They may increase over time and can affect appearance; treatment is chosen according to number, size, and location.

\medterm{Familial Paroxysmal Peritonitis} An inherited inflammatory disorder causing repeated short attacks of fever with severe abdominal, chest, or joint pain. Without treatment, ongoing inflammation can damage the kidneys; preventive treatment can reduce attacks and complications.

\synonyms
Familial Mediterranean fever

\medterm{Familial Testotoxicosis} An inherited condition in which boys produce excess testosterone without normal brain stimulation. It causes early sexual development, rapid growth, and advanced bone age but may reduce adult height without treatment.

\medterm{Family} People related by biology, law, or social bonds. Family history, shared genes, environment, and support relationships can affect health risks, decisions, caregiving, and recovery.

\medterm{Family Involvement} Participation of family members or chosen support people in healthcare planning, treatment, rehabilitation, or daily care, with the patient’s permission and attention to privacy, safety, preferences, and caregiver burden.

\medterm{Family Planning} Personal and healthcare decisions about whether and when to have children. It includes contraception, fertility care, preconception health, pregnancy planning, and support for informed reproductive choices.

\medterm{Family Study} A structured review of health information in relatives to understand how a disease or trait occurs in a family. It may include family history, medical records, examination, genetic counseling, or testing.

\medterm{Family Therapy} Counseling involving family members or other close relationships to improve communication, manage conflict, and support recovery. It may help with mental illness, addiction, chronic disease, caregiving stress, or family adjustment.


\medterm{Famophos} An organophosphate insecticide that can poison the nervous system by blocking acetylcholinesterase. Exposure may cause sweating, salivation, vomiting, pinpoint pupils, muscle weakness, breathing failure, or death and requires emergency care.

\medterm{Famotidine} A medicine that blocks histamine signals to reduce stomach-acid production. It treats reflux, heartburn, and ulcers; dose reduction may be needed in kidney disease, and headache, diarrhea, or confusion can occur.

\medterm{Fanconi Anemia} An inherited disorder affecting DNA repair and bone-marrow blood-cell production. It may cause low blood counts, infections, bleeding, short stature, birth differences, organ problems, and a high risk of leukemia or other cancers.

\medterm{Fanconi Syndrome} A kidney-tubule disorder causing excess loss of glucose, phosphate, bicarbonate, amino acids, and other substances in urine. It can lead to dehydration, abnormal blood acidity, weak bones, poor growth, or kidney disease.

\medterm{FAP} An inherited condition causing many precancerous polyps in the colon and rectum. Colorectal cancer is very likely without preventive treatment, so regular specialist surveillance and risk-reducing surgery are often needed.

\synonyms
Familial adenomatous polyposis

\medterm{Farad} The standard unit of electrical capacitance, equal to one coulomb per volt. It is used in medical equipment and electrophysiology to describe how much electrical charge a component can store.

\medterm{Farber Lipogranulomatosis} A rare inherited storage disorder caused by deficient acid ceramidase. Fatty material accumulates and may cause painful skin lumps, joint stiffness, hoarseness, bone changes, breathing problems, and neurologic disease.

\medterm{Farmer'S Lung} Lung inflammation caused by repeatedly inhaling moldy hay, grain, or other farm dust. It can cause fever, cough, chills, and breathlessness hours after exposure and may lead to permanent scarring if exposure continues.

\medterm{Farmer’S Lung} An immune lung disease caused by inhaling particles from moldy hay, grain, or other farm materials. Repeated exposure can cause cough, fever, and breathlessness and eventually permanent lung scarring; avoiding the source is central to care.

\medterm{Farsightedness} A focusing problem in which close objects look blurred because light would focus behind the retina. It may cause eyestrain or headaches and is usually corrected with glasses or contact lenses.

\synonyms
Hyperopia

\medterm{Fart} A colloquial term for passing intestinal gas through the anus. Gas comes from swallowed air and bacterial breakdown of food; frequent painful or newly changed gas may reflect diet, constipation, intolerance, or digestive disease.

\medterm{Fascia} Sheets or bands of connective tissue that surround, separate, support, and connect muscles, organs, nerves, and blood vessels. Injury or inflammation of fascia can cause pain, stiffness, or restricted movement.

\medterm{Fascia Adherens} A strong connection joining neighboring cells, especially heart-muscle cells. It links the internal contractile framework of one cell to the membrane of another so force can pass through the tissue.

\medterm{Fascial Loop} A loop or sling made from a person’s tissue or synthetic material to support a body structure. For example, a sling may support the urethra in selected treatment of stress urinary incontinence.

\medterm{Fasciculation} A brief, involuntary twitch visible under the skin from contraction of a small group of muscle fibers. It may be harmless, but widespread or persistent twitching with weakness needs neurologic evaluation.

\medterm{Fasciitis} Inflammation of the connective tissue around muscles and other structures. It may be localized, as in plantar fasciitis, or rapidly destructive, as in necrotizing infection requiring emergency surgery and antibiotics.

\medterm{Fasciola} A genus of parasitic flatworms, or liver flukes, that infect some animals and people. Infection is acquired by swallowing immature parasites on contaminated water plants or water and may affect the liver and bile ducts.

\medterm{Fasciola Hepatica} A liver fluke that infects people after contaminated aquatic plants or water are swallowed. Immature worms migrate through the liver, while adult worms live in bile ducts and can cause pain, fever, or obstruction.

\medterm{Fascioliasis} Infection with liver flukes after eating contaminated aquatic plants or drinking contaminated water. It may cause fever, right-upper-abdominal pain, enlarged liver, or high eosinophils and later cause bile-duct pain, infection, or blockage.

\medterm{Fasciolopsis} A genus of large intestinal flukes that can infect people who eat contaminated aquatic plants. Heavy infection may cause abdominal pain, diarrhea, swelling, poor nutrition, or intestinal obstruction.

\medterm{Fasciolopsis Buski} A large intestinal fluke acquired by eating aquatic plants carrying immature parasites. Heavy infection may cause abdominal pain, diarrhea, intestinal inflammation, malnutrition, or swelling.

\medterm{Fasciotomy} Emergency surgery that cuts tight fascia to release dangerous pressure in a muscle compartment. It may prevent permanent nerve and muscle damage after crush injury, fracture, burns, or loss and return of blood flow.


\medterm{Fast Heart Rate} A heart rate faster than expected for a person’s age and activity. Exercise, fever, pain, stress, dehydration, anemia, medicines, or an abnormal rhythm may cause it; chest pain or fainting is urgent.

\medterm{Fast Heart Rate - Tachycardia} A rapid heartbeat that may be a normal response to exercise, fever, pain, anxiety, dehydration, anemia, or medicines, or may reflect an abnormal rhythm. Chest pain, fainting, or severe breathlessness requires urgent care.

\medterm{Fast-Twitch Fiber} A muscle fiber capable of quick, powerful contraction but more rapid fatigue than a slow-twitch fiber. It supports sprinting, jumping, lifting, and other brief high-intensity activity.

\synonyms
Type II muscle fiber

\medterm{Fasting} Going without food, and sometimes without drinks, for a defined period. It may be required before blood tests, anesthesia, surgery, or certain procedures, or undertaken for personal or religious reasons.

\medterm{Fasting Blood Glucose} Blood sugar measured after no caloric intake, usually for at least eight hours. It helps screen for and diagnose diabetes or impaired glucose regulation, but results must be interpreted with the clinical context.

\medterm{Fasting Guidelines} Instructions about when to stop food and drinks before anesthesia, surgery, or a test. Timing depends on the procedure, age, health, and local policy; the treating team’s instructions must take priority.

\medterm{Fasting Lipid Profile} A blood test measuring cholesterol and triglycerides after a period without food. It helps assess cardiovascular risk and guide treatment; fasting is particularly useful when triglycerides need accurate measurement.

\synonyms
Lipid panel

\medterm{Fasting Plasma Glucose Test} A blood test measuring blood glucose after fasting, usually at least eight hours. It helps identify diabetes or prediabetes and is interpreted with repeat testing or other results when needed.

\medterm{Fat} Energy-rich substances that provide fuel, insulation, and components for cell membranes and hormones. Dietary fats also help absorb vitamins A, D, E, and K; types include saturated, unsaturated, and trans fats.

\synonyms
Lipid

\medterm{Fat Embolism} Entry of fat droplets into the bloodstream, usually after major trauma, a thigh or pelvic fracture, or orthopedic surgery. It may cause breathing difficulty, confusion, rash, or low oxygen and can become life-threatening.

\medterm{Fat Measurement} Estimation of body-fat amount or distribution using methods such as skinfolds, bioimpedance, DEXA, imaging, or laboratory analysis. Results vary by method and should be interpreted with overall health, strength, and metabolic findings.

\medterm{Fat Necrosis} Death of fat cells, often after pancreatitis, injury, or breast surgery. It can cause pain, inflammation, or a firm lump that resembles cancer on examination or imaging and may require testing to distinguish the cause.

\medterm{Fat Pad} A localized cushion of fatty tissue that protects, supports, insulates, or permits movement around an organ or joint. Inflammation or injury of a fat pad can cause localized pain and swelling.

\medterm{Fat-Soluble Toxins} Poisonous chemicals that dissolve in body fat and may remain for a long time. Some pesticides, industrial chemicals, polychlorinated biphenyls, and dioxins can accumulate after repeated exposure and harm several organs.

\medterm{Fat-Soluble Vitamin} A vitamin that dissolves in fat and can be stored in body tissues. Vitamins A, D, E, and K are fat-soluble; deficiency or excessive supplementation can cause illness.

\medterm{Fat-Soluble Vitamins} Vitamins A, D, E, and K, which dissolve in fat and can accumulate in the body. They support vision, bones, immunity, cell protection, and blood clotting; excessive supplements can be toxic.

\medterm{Fatal Familial Insomnia} A very rare inherited prion disease that progressively disrupts sleep, blood-pressure and temperature control, movement, and thinking. It causes severe neurologic decline and is ultimately fatal; specialist genetic and supportive care is needed.

\medterm{Fatigue} Persistent tiredness or reduced energy that limits physical or mental activity and may not improve fully with rest. Sleep problems, anemia, infection, hormone or organ disease, medicines, mood disorders, and neurologic conditions can cause it.

\medterm{Fatigue (Cancer)} Persistent tiredness or low energy related to cancer or its treatment. Anemia, pain, poor sleep, infection, medicines, emotional distress, and the cancer itself may contribute; evaluation can identify treatable causes.

\synonyms
Cancer-related fatigue

\medterm{Fatty Acid} A building block of fats that can provide energy, form cell membranes, and serve as a starting material for signaling chemicals. Fatty acids may be saturated, monounsaturated, or polyunsaturated.

\medterm{Fatty Acid Synthase} A group of linked enzymes that helps cells make fatty acids, especially palmitate. These fatty acids are used for energy storage, cell membranes, and other cellular functions.
\medterm{Fatty Acids} Components of fats used for energy, cell membranes, and chemical signals. Their effects on health depend on their type and amount.
\medterm{Fatty Degeneration} Abnormal buildup of fat inside cells, especially liver cells, caused by metabolic problems, alcohol, toxins, poor nutrition, obesity, or other injury. The change may reverse or progress to inflammation and scarring.

\medterm{Fatty Liver (Steatosis)} Excess fat stored in liver cells. Alcohol, obesity, diabetes, some medicines, and other causes may contribute; it may remain harmless or progress to liver inflammation, scarring, cirrhosis, or cancer.
\synonyms
Hepatic steatosis; fatty liver

\medterm{Fatty Liver} Alternate index label; see \textit{Fatty Liver (Steatosis)}.
\medterm{Fatty Liver Disease} Excess fat in the liver, often associated with obesity, diabetes, abnormal cholesterol, or alcohol use. Some people develop inflammation and scarring, so cause, metabolic health, and liver risk should be assessed.

\synonyms
Hepatic steatosis; steatotic liver disease

\medterm{Fatty Streak} An early artery-wall change in which cholesterol-rich immune cells collect beneath the lining. It can be an early stage of atherosclerosis, although some fatty streaks never progress to a harmful plaque.

\medterm{Favism} Rapid destruction of red blood cells after eating fava beans, usually in a person with glucose-6-phosphate dehydrogenase deficiency. It may cause jaundice, dark urine, weakness, breathlessness, or severe anemia and needs medical care.

\medterm{Favus} A long-lasting fungal infection of the scalp, usually causing yellow cup-shaped crusts around hairs. Without treatment it can destroy follicles and cause permanent scarring hair loss.

\medterm{Fc Receptor} A protein on an immune cell that binds the tail end of an antibody. This helps the cell recognize, ingest, or destroy the antibody’s target and can also influence inflammation.

\medterm{Fear} An emotional and body response to a perceived threat that prepares a person to protect themselves. Persistent or excessive fear can cause avoidance, panic, sleep problems, and impaired daily functioning.

\medterm{Fears} Feelings of threat or apprehension. Fears are normal when proportionate to danger, but persistent, excessive, or disabling fears may indicate an anxiety disorder and can often be treated.

\medterm{Febrile} Relating to fever, a temporary rise in body temperature usually caused by infection, inflammation, or another illness. Fever differs from heat-related hyperthermia; severe illness, dehydration, confusion, or fever with very low white cells is urgent.

\medterm{Febrile (Warm) And Cold Agglutinins} Antibodies that make red blood cells clump at warmer or colder temperatures. They can interfere with laboratory testing or cause autoimmune destruction of red cells and anemia.

\medterm{Febrile Convulsion} A seizure during fever in a young child without evidence of brain infection or another immediate cause. Most are brief and generalized, but a prolonged, repeated, or one-sided seizure requires emergency assessment.

\medterm{Febrile Convulsions} Seizures triggered by fever in young children, usually between six months and five years, without evidence of infection or another direct brain problem. Most resolve quickly, but first or prolonged seizures require medical assessment.

\medterm{Febrile Neutropenia} Fever in a person with a dangerously low neutrophil count, the white cells that fight infection. Infection may be present without obvious symptoms, so urgent examination, cultures, and prompt antibiotics are often needed.

\medterm{Febrile Seizure} A seizure during fever in a child, usually between six months and five years, without evidence of brain infection or a previous seizure without fever. Most are brief, but prolonged, repeated, or one-sided events require urgent assessment.

\synonyms
Seizure, Febrile

\medterm{Febrile Seizures} Seizures associated with fever in young children without evidence of brain infection or a previous unprovoked seizure. Most do not cause lasting harm, but the first seizure or any prolonged event needs medical evaluation.

\synonyms
Seizure Febrile (Febrile Seizures)
Seizure Fever Induced (Febrile Seizures)


\medterm{Febuxostat} A medicine that lowers uric-acid production for long-term prevention of gout attacks. It does not relieve sudden gout pain immediately and may require monitoring for liver problems and cardiovascular risk.

\medterm{FEC} Alternate index label; see \textit{FECD}.

\medterm{Fec Protocol} A chemotherapy regimen combining fluorouracil, epirubicin, and cyclophosphamide for selected cancers. It can cause low blood counts, infection, nausea, hair loss, mouth sores, infertility, and heart or bladder-related toxicity.

\medterm{Fecal Calprotectin} A stool test measuring a protein released during intestinal inflammation. A raised result can support evaluation for inflammatory bowel disease, but infection, medicines, age, and other conditions can also raise it.

\medterm{Fecal Culture} Laboratory testing that grows bacteria from a stool sample to investigate infectious diarrhea. It may identify a treatable cause and guide antibiotics, but some infections require toxin, molecular, or parasite testing instead.

\medterm{Fecal Elastase} A stool test estimating whether the pancreas releases enough digestive enzymes. A low result may suggest pancreatic insufficiency, causing poor digestion, greasy stools, weight loss, or vitamin deficiency.

\medterm{Fecal Fat} Measurement of excess fat in stool. High fecal fat suggests poor digestion or absorption from pancreatic disease, reduced bile, intestinal disease, or another condition and may explain greasy, difficult-to-flush stools.

\medterm{Fecal Immunochemical Test} A stool screening test that detects human blood, mainly from the lower intestine, using antibodies. A positive result does not diagnose cancer but usually requires colon evaluation, often colonoscopy.

\medterm{Fecal Impaction} A large, hard mass of stool trapped in the rectum or colon. It may cause pain, abdominal swelling, nausea, urinary symptoms, bowel blockage, or watery leakage around the mass and may require professional removal.

\medterm{Fecal Incontinence} Inability to control passage of stool or gas, ranging from soiling to complete loss of bowel control. Causes include diarrhea, constipation, sphincter injury, nerve disease, rectal prolapse, and cognitive impairment.

\synonyms
Incontinence, Bowel
Incontinence, Fecal

\medterm{Fecal Microbiota Transplant} Transfer of carefully screened and prepared donor stool or its microbial contents into the intestine to restore helpful microbes, mainly for selected recurrent C. difficile infection. Infection transmission remains a possible risk.

\medterm{Fecal Microbiota Transplantation} Transfer of screened donor stool or microbial material into a patient’s intestine to restore microbial balance. It is used mainly for selected recurrent C. difficile infections after standard therapy and requires regulated donor screening.

\medterm{Fecal Occult Blood Test} A stool test that detects blood too small to see. It may screen for colorectal cancer or investigate digestive-tract bleeding, but it does not identify the cause; a positive result usually requires further colon evaluation.

\medterm{Fecal Smear} Examination of a thin stool sample with a microscope or stain. Depending on the test, it may help identify parasites, inflammatory cells, blood, fat, or other abnormalities.

\medterm{Fecalith} A hard mass of dried stool, often in the appendix or colon. It may contribute to appendicitis, blockage, inflammation, or perforation.

\medterm{Fecaluria} Stool, gas, or a fecal smell in urine, usually indicating an abnormal connection between the intestine and urinary tract. It often accompanies recurrent urinary infection and requires medical evaluation.

\medterm{FECD} An abbreviation for Fuchs endothelial corneal dystrophy, a progressive loss of corneal cells that normally remove fluid. It can cause morning blur, glare, painful swelling, and reduced vision; advanced disease may need a corneal transplant.

\synonyms
Fuchs dystrophy

\medterm{Feces} Solid or semisolid waste passed from the bowel, containing water, undigested material, bacteria, shed cells, and metabolic waste. Changes in stool can provide clues about digestion, bleeding, infection, or inflammation.

\medterm{Fecolith} A hard mass of dried stool that can block the appendix or bowel and contribute to inflammation, appendicitis, or perforation.

\medterm{Fecundity} The biological ability to conceive and produce a live birth. It depends on age, ovulation, ovarian reserve, fallopian tubes, uterus, sperm health, and other factors; it differs from the chance of pregnancy in one cycle.

\medterm{Feedback} Information about a system’s output that changes later activity. Negative feedback stabilizes functions such as temperature, glucose, and hormones, while positive feedback strengthens processes such as labor contractions.

\medterm{Feedback Inhibition} A control process in which the end product of a chemical pathway slows an earlier step. This prevents unnecessary buildup of the product and helps cells use nutrients and energy efficiently.

\medterm{Feedback Mechanism} A regulatory process in which a system’s result influences what happens next. Negative feedback maintains stability, while positive feedback amplifies a response until a defined endpoint is reached.

\medterm{Feeding Jejunostomy} Placement of a feeding tube into the jejunum, part of the small intestine, when eating or stomach feeding is unsafe or impossible. It provides nutrition but can cause infection, blockage, leakage, or bowel injury.

\medterm{Feeding Methods} Ways of providing food or liquid, including breast or bottle feeding, oral feeding, tube feeding, and intravenous nutrition. The safest method depends on swallowing ability, age, illness, and nutritional needs.

\medterm{Feeding Patterns And Diet - Babies And Infants} Guidance on how babies receive breast milk, formula, and later complementary foods. Advice should match age, growth, development, swallowing safety, allergies, and medical needs.

\medterm{Feeding Patterns And Diet - Children 6 Months To 2 Years} Guidance on introducing varied foods and drinks while continuing appropriate milk feeding. Portion size, texture, allergy prevention, choking safety, growth, and iron intake are important.

\medterm{Feeding Poor} Taking too little food or fluid for age or health needs. In an infant it may signal illness, swallowing difficulty, breathing problems, pain, reflux, or developmental difficulty and requires assessment if persistent.

\medterm{Feeding Tube} A flexible tube used to deliver nutrition, fluids, or medicines to the stomach or intestine when eating is inadequate or unsafe. Risks include blockage, leakage, infection, aspiration, and injury.

\medterm{Feeding Tube - Infants} A tube delivering breast milk, formula, fluids, or medicines to an infant unable to feed safely or adequately by mouth. The route, feeding plan, and duration depend on swallowing, breathing, growth, and the underlying illness.

\medterm{Feeding Tube Insertion - Gastrostomy} Placement of a tube through the abdominal wall into the stomach to provide nutrition, fluids, or medicines. It may be temporary or long-term and can cause bleeding, infection, leakage, blockage, or injury.

\medterm{Feingold Diet} A restrictive diet that avoids certain food additives and some naturally occurring chemicals, promoted for behavior or attention problems. Evidence of benefit is limited, and restriction can cause nutritional problems without professional guidance.

\medterm{Felbamate} An antiseizure medicine reserved for selected severe or drug-resistant epilepsy. It can rarely cause life-threatening bone-marrow failure or liver failure, so specialist prescribing and regular blood and liver tests are essential.

\medterm{Felodipine} A calcium-channel blocker that relaxes artery walls and lowers blood pressure. It may cause ankle swelling, flushing, headache, fast heartbeat, or gum enlargement and should be taken as prescribed.

\medterm{Felon} A painful bacterial infection in the fleshy fingertip, often after a small puncture. Pus and pressure can damage tissue, bone, or tendons, so prompt assessment and sometimes drainage and antibiotics are needed.

\medterm{Felty Syndrome} A rare complication of long-standing rheumatoid arthritis involving a very low neutrophil count and often an enlarged spleen. It increases infection risk and requires specialist treatment of arthritis and neutropenia.

\medterm{Felty'S Syndrome} A rare complication of long-standing rheumatoid arthritis causing neutropenia and often an enlarged spleen. It raises the risk of serious infection and may require disease-modifying treatment, growth factors, or other specialist care.

\medterm{Female} Describing the sex category generally associated with ovaries, eggs, and female reproductive anatomy. Biological traits vary among individuals, and sex, gender identity, and gender expression are distinct concepts.

\medterm{Female Athlete Triad} A pattern of low energy availability, menstrual disturbance, and reduced bone strength in physically active females. Inadequate nutrition or excessive exercise can increase the risk of stress fractures and other health problems.

\medterm{Female Breast} A mammary gland and its skin, fat, connective tissue, ducts, and lobules on the front of the chest. Its structure changes with age, hormones, pregnancy, breastfeeding, and body composition.

\medterm{Female Circumcision} Removal or injury of part of the external female genitalia for nonmedical reasons. It has no health benefit and can cause severe pain, bleeding, infection, urinary and sexual problems, childbirth complications, and lasting psychological harm.

\medterm{Female Condom} A soft internal sheath placed in the vagina before sex that forms a barrier against semen and can reduce pregnancy and transmission of some sexually transmitted infections when used correctly.

\medterm{Female Condoms} Flexible internal barriers placed in the vagina before sex to help prevent pregnancy and reduce exposure to semen and some sexually transmitted infections. They should not be used at the same time as an external condom.

\medterm{Female Genital Mutilation} Cutting or removing external female genital tissue for nonmedical reasons. It has no health benefit and can cause pain, bleeding, infection, urinary and sexual problems, childbirth complications, and long-term psychological harm.

\medterm{Female Genitalia} The external and internal reproductive organs, including the vulva, vagina, cervix, uterus, fallopian tubes, and ovaries. These structures support urination, sexual function, menstruation, fertility, pregnancy, and birth.

\medterm{Female Gonad} An ovary, the reproductive gland that stores and releases eggs and produces hormones including estrogen and progesterone.

\medterm{Female Infertility} Difficulty achieving pregnancy because of problems with ovulation, ovarian reserve, fallopian tubes, uterus, hormones, or other factors. Evaluation considers both partners and may be recommended after a defined period of trying or sooner with risk factors.

\synonyms
Infertility, Female

\medterm{Female Orgasmic Disorder (Anorgasmia)} Persistent or recurrent difficulty achieving orgasm that causes personal distress. It may be lifelong or acquired and can relate to health conditions, medicines, mood, trauma, sexual pain, stimulation, or relationship factors.

\synonyms
Anorgasmia

\medterm{Female Pattern Baldness} Gradual scalp hair thinning, usually widest near the part or crown, influenced by genes and hormones. Iron deficiency, thyroid disease, medicines, and other causes should be considered when loss is new or rapid.

\medterm{Female Pelvis} The pelvic bones and soft tissues in a person with typical female anatomy. It is often broader with a wider outlet than typical male anatomy, but substantial individual variation exists.

\medterm{Female Phenotype} The observable anatomical, physiologic, and developmental traits associated with female sex. These traits vary with genes, hormones, development, age, and individual biology.

\medterm{Female Reproductive System} The ovaries, fallopian tubes, uterus, cervix, vagina, vulva, and related hormones. It supports egg release, menstruation, pregnancy, birth, and sexual function and can be affected by developmental, hormonal, infectious, structural, or cancerous disorders.

\medterm{Female Sexual Dysfunction} Persistent sexual difficulty causing distress, such as low desire, arousal or lubrication problems, difficulty with orgasm, or pain. Medical conditions, medicines, mood, trauma, relationships, and life changes may contribute.

\synonyms
Sexual Dysfunction, Female

\medterm{Female Sexual Interest/Arousal Disorder} Persistent or recurrent low sexual interest or arousal causing personal distress. Medical conditions, medicines, mood, trauma, pain, relationship factors, hormonal changes, and other social or psychological factors may contribute.

\medterm{Female Urethral Opening} The opening through which urine leaves the female urethra, located in the vulvar vestibule in front of the vaginal opening. Its position and appearance may vary normally.

\medterm{Feminization} Development of physical or hormonal traits commonly associated with female puberty, such as breast growth or reduced body hair. It may be normal, part of gender-affirming treatment, or caused by hormones, medicines, or disease.

\medterm{Femoral} Relating to the femur or thigh, including the thigh bone, artery, vein, nerve, or neck of the femur.

\medterm{Femoral Artery} A major artery carrying blood through the groin and thigh to the leg. Its pulse can be felt in the groin, and it is used for some catheter procedures; injury or blockage can threaten the limb.

\medterm{Femoral Hernia} A bulge of abdominal tissue through an opening below the groin ligament and beside the femoral vessels. It has a relatively high risk of trapping or losing blood supply to the contents and often needs surgical assessment.

\medterm{Femoral Hernia Repair} Surgery that returns trapped tissue to the abdomen and closes or reinforces the opening in the groin. Repair is often recommended because femoral hernias can become trapped or strangulated.

\medterm{Femoral Nerve} A large nerve from the lower back that supplies muscles that flex the hip and straighten the knee and provides sensation to the front thigh and inner leg. Injury can cause weakness, numbness, or falls.

\medterm{Femoral Nerve Dysfunction} Reduced function of the nerve supplying the front thigh and knee-straightening muscles. It can cause hip or knee weakness and altered sensation and may result from diabetes, compression, injury, surgery, or pelvic disease.

\medterm{Femoral Neuropathy} Damage to the femoral nerve causing weakness in hip flexion or knee extension and numbness or pain in the front thigh or inner leg. Causes include diabetes, injury, bleeding, compression, and pelvic disease.
\medterm{Femoral Pulse} The pulse felt over the femoral artery in the groin. Comparing it with pulses elsewhere helps assess blood flow to the leg and can provide clues about arterial blockage or shock.

\medterm{Femoral Shaft Fractures} Breaks in the middle portion of the thigh bone, often from high-energy trauma but sometimes from weakened bone. They can cause major blood loss, deformity, and inability to walk and usually require urgent stabilization and specialist treatment.

\medterm{Femoral Triangle} A triangular area at the upper front of the thigh bordered by the groin ligament and nearby muscles. It contains the femoral nerve, artery, vein, and lymph nodes and is an important examination and access site.

\medterm{Femoral Vein} The major deep vein of the thigh, carrying blood from the leg toward the pelvis and heart. A clot in it can cause deep-vein thrombosis and pulmonary embolism.

\medterm{Femtogram} A unit of mass equal to one quadrillionth of a gram, or $10^{-15}$ gram. It is used in laboratory measurements of extremely small quantities.

\medterm{Femtoliter} A unit of volume equal to one quadrillionth of a liter, or $10^{-15}$ liter. It is used to describe very small laboratory or cellular volumes.

\medterm{Femur} The thigh bone and the longest, strongest bone in the body. It joins the hip to the knee and includes a head, neck, shaft, and lower end; fractures can threaten mobility and blood supply.

\medterm{Femur Fracture} A break in the thigh bone caused by trauma or weakened bone. It may cause severe pain, deformity, blood loss, and inability to bear weight; treatment depends on location, alignment, and overall health.


\medterm{Femur Length} The ultrasound-measured length of a fetus’s thigh bone. It helps estimate pregnancy age and assess growth, but interpretation depends on other measurements, family size, pregnancy dating, and possible skeletal or chromosomal conditions.

\medterm{Fencer'S Pose} A posture in which turning the head makes the arm and leg on that side straighten while the opposite arm and leg bend. It is a normal reflex in young infants but may indicate a brain or nerve problem if it persists or appears unexpectedly.

\medterm{Fenfluramine} Fenfluramine is a serotonergic medicine used in some jurisdictions for selected severe epileptic encephalopathies; historical appetite-suppressant use was limited by valvular heart disease and pulmonary hypertension, so cardiac monitoring is required.

\medterm{Fenitrothion} A poisonous organophosphate insecticide. Exposure can cause sweating, excess saliva, vomiting, diarrhea, small pupils, muscle twitching, breathing difficulty, seizures, or weakness because it disrupts nerve signals; suspected poisoning needs urgent medical help.

\medterm{FeNO (Fractional Exhaled Nitric Oxide)} A breath test that estimates type 2, often eosinophilic, airway inflammation. The result is interpreted with symptoms, age, treatment, and other clinical findings.

\synonyms
Fractional Exhaled Nitric Oxide

\medterm{Fenobucarb} A carbamate insecticide that can cause nerve poisoning after significant absorption, producing sweating, vomiting, small pupils, muscle weakness, breathing difficulty, or seizures.

\medterm{Fenofibrate} A medicine that lowers blood triglycerides and can modestly change cholesterol levels by increasing the liver's processing of fats. It is used for selected lipid disorders; liver, kidney, muscle, and gallbladder risks require monitoring.

\medterm{Fenoldopam} A short-acting dopamine D1-receptor agonist used intravenously for severe hypertension in selected settings; it can cause hypotension, flushing, headache, and increased intraocular pressure.

\medterm{Fenoprofen} Fenoprofen is a nonsteroidal anti-inflammatory drug used for pain and inflammatory arthritis; gastrointestinal bleeding, renal injury, fluid retention, and cardiovascular events are important risks.

\medterm{Fenoprofen Calcium Overdose} Taking too much fenoprofen, an anti-inflammatory pain medicine, can cause nausea, vomiting, stomach bleeding, drowsiness, kidney injury, breathing or acid-base problems, seizures, or coma. Suspected overdose requires immediate poison-control or emergency advice; do not induce vomiting.

\medterm{Fenoterol} A short-acting beta-2 bronchodilator that relaxes airway muscle and relieves bronchospasm. It may cause tremor, nervousness, low potassium, or a fast or irregular heartbeat and is not approved everywhere.

\medterm{Fenoxycarb} An insect-growth regulator that imitates juvenile hormone and disrupts development of immature insects. It is used for pest control, and exposure should be minimized because products can irritate skin, eyes, or airways.

\medterm{Fentanyl Citrate} The citrate salt of fentanyl, a potent opioid analgesic used for severe pain and anesthesia; respiratory depression, sedation, dependence, and overdose are major risks.

\medterm{Fentanyl Derivatives} Strong synthetic opioid medicines related to fentanyl, including sufentanil, remifentanil, and alfentanil. They may relieve severe pain or support anesthesia but can rapidly slow breathing, cause overdose, and produce dependence, so close monitoring is essential.

\medterm{Fenthion} A poisonous organophosphate insecticide. Exposure can cause sweating, excess saliva, vomiting, diarrhea, small pupils, muscle twitching, breathing difficulty, seizures, or weakness because it disrupts nerve signals; suspected poisoning needs urgent medical help.

\medterm{Fermentation} A metabolic process in which microorganisms or cells convert organic substrates without using oxidative respiration, producing products such as acids, alcohols, or gases.

\medterm{Ferritin} A protein that stores iron inside cells and releases it when needed. Blood ferritin often reflects body iron stores, but inflammation, infection, liver disease, or cancer can raise it even when iron is low.

\medterm{Ferritin Blood Test} A blood test measuring stored iron indirectly. Low ferritin strongly supports iron deficiency; high ferritin may reflect iron excess, inflammation, infection, liver disease, or cancer and must be interpreted with other tests.

\medterm{Ferroptosis} A controlled form of cell death caused by iron-dependent damage to cell fats and failure of antioxidant protection. It differs from apoptosis and ordinary tissue necrosis and is being studied in cancer and other diseases.

\medterm{Ferrous Sulfate} An oral iron supplement used to prevent or treat iron-deficiency anemia. It may cause nausea, constipation, abdominal discomfort, or dark stools; accidental overdose can be life-threatening, especially in children.

\medterm{Fertile Period} The several days in a menstrual cycle when sex can result in pregnancy, including the days before and shortly after ovulation. Timing varies, and calendar estimates are not reliable contraception for everyone.

\medterm{Fertilisation} Union of a sperm and an egg to form a zygote, the first cell of an embryo. It usually occurs in a fallopian tube before the developing embryo moves toward the uterus.

\medterm{Fertility} The biological ability to conceive and produce offspring. It depends on ovulation, sperm production, reproductive anatomy, hormones, age, general health, and other factors affecting either partner.

\medterm{Fertility Domain} The part of health concerned with reproductive function, ability to conceive, fertility goals, infertility, preservation of fertility, and related medical, emotional, and social care.

\medterm{Fertility Drugs} Medicines used to stimulate ovulation, support reproductive treatment, or improve the chance of pregnancy. They can cause side effects and multiple pregnancy, so selection and monitoring require fertility-specialist care.

\medterm{Fertility Preservation} Steps taken to protect the possibility of future pregnancy before treatment or disease may damage reproductive function. Options include freezing eggs, embryos, sperm, or selected reproductive tissue, depending on age, timing, and goals.

\medterm{Fertilization} Fusion of sperm and egg genetic material to form a zygote, usually in a fallopian tube. This begins embryonic development, which may continue after the zygote reaches the uterus and implants.

\medterm{Ferumoxytol} An intravenous iron preparation used for selected adults with iron-deficiency anemia. It can cause allergic reactions, low blood pressure, or iron overload, so monitoring is required during and after infusion.

\medterm{FESS} Functional endoscopic sinus surgery uses instruments passed through the nostrils to open blocked sinus drainage pathways and remove selected diseased tissue. It may improve chronic sinus disease but can cause bleeding, infection, eye injury, or fluid leakage.

\medterm{Festinating Gait} Walking with increasingly short, quick steps and a forward-leaning posture, commonly associated with Parkinson disease or other movement disorders. It increases fall risk and may improve with specialist treatment and physical therapy.

\medterm{Fetal Acid-Base Status} An assessment of a fetus’s oxygen supply and blood acidity, often using umbilical-cord blood gases at birth. Results must be interpreted with the sampling site, timing, fetal condition, and complete clinical record.

\medterm{Fetal Alcohol Spectrum Disorder} Lifelong physical, learning, behavioral, or developmental effects caused by alcohol exposure before birth. Effects vary widely and may include growth problems, attention difficulties, poor coordination, or characteristic facial features.

\synonyms
Fetal Alcohol Spectrum Disorders


\medterm{Fetal Alcohol Syndrome} The most severe recognized form of fetal alcohol spectrum disorder, involving characteristic facial features, poor growth, and brain-development problems after prenatal alcohol exposure. No known safe amount of alcohol exists during pregnancy.

\synonyms
Alcohol Pregnancy (Fetal Alcohol Syndrome)

\medterm{Fetal Attitude} The position of a fetus’s head, arms, legs, and trunk in relation to one another. Normal flexion helps birth; unusual extension or flexion can change the presenting part and affect labor.

\medterm{Fetal Bradyarrhythmia} An unusually slow or irregular fetal heart rhythm. Causes include conduction block, extra beats, maternal or fetal disease, or medicines; persistent abnormalities require prompt specialist assessment.

\medterm{Fetal Bradycardia} A fetal heart-rate baseline below 110 beats per minute for at least ten minutes. Medicines, maternal illness, fetal heart disease, low oxygen, cord problems, or placental separation may cause it and persistent cases require urgent obstetric assessment.

\medterm{Fetal Circulation} Blood flow before birth in which the placenta supplies oxygen and nutrients and special connections direct much blood around the lungs, which are not yet used for breathing. These connections normally close after birth.

\medterm{Fetal Death} Death of a fetus before delivery. Medical and legal definitions vary by gestational age and location; care includes confirming the diagnosis, supporting the family, and discussing delivery and follow-up.

\medterm{Fetal Development} Growth and maturation from the end of the embryonic period until birth. Organs enlarge and mature, movements develop, and body systems gradually acquire the functions needed after delivery.

\medterm{Fetal Disease} A disease, structural abnormality, or developmental problem affecting a fetus before birth. Causes include genetic changes, infection, medicines, toxins, and poor placental function; ultrasound and prenatal tests may identify it.

\medterm{Fetal Distress} A nonspecific term for concerning fetal findings during labor, especially abnormal heart-rate patterns that may indicate low oxygen, cord compression, infection, or placental problems. Clinicians assess the whole pattern and may act urgently.

\medterm{Fetal Echocardiography} A specialized ultrasound examination of the fetal heart used to assess cardiac structure, rhythm, and function, including possible congenital heart disease.

\medterm{Fetal Growth Restriction} A fetus that is smaller than expected because it may not be reaching its growth potential, often because the placenta is not supplying enough oxygen or nutrients. Repeat ultrasound, fluid, blood-flow, and well-being checks guide care.

\medterm{Fetal Heart} The developing heart of a fetus, whose rate and rhythm can be assessed by ultrasound or electronic monitoring to evaluate fetal well-being.

\medterm{Fetal Heart Rate} The number of times a fetus’s heart beats each minute, usually about 110–160 during pregnancy and labor. It can be checked with a Doppler or electronic monitor; changes may reflect normal activity, medicines, or reduced oxygen.

\medterm{Fetal Hemoglobin} The main oxygen-carrying protein in a fetus’s red blood cells. It binds oxygen more strongly than adult hemoglobin, helping the fetus obtain oxygen from the placenta; its amount normally falls after birth.

\medterm{Fetal Hydrops} Severe fluid buildup in at least two parts of a fetus, such as under the skin, around the lungs or heart, or in the abdomen. Causes include anemia, heart disease, infection, chromosome conditions, and blood-group incompatibility and require specialist care.

\medterm{Fetal Hypoxia} Too little oxygen reaching a fetus before or during birth. Placental disease, maternal breathing or circulation problems, cord compression, or prolonged labor can contribute; abnormal heart-rate patterns require prompt obstetric assessment.

\medterm{Fetal Kick Counts (Fetal Movement Counting)} Observation and recording of fetal movements by a pregnant person. A noticeable reduction from the usual pattern should prompt timely obstetric advice.

\synonyms
Fetal Movement Counting

\medterm{Fetal Lie} The direction of a fetus’s long body axis compared with the pregnant person’s spine. The fetus may lie lengthwise, sideways, or diagonally; a sideways or diagonal position near birth can affect delivery planning.

\medterm{Fetal Macrosomia} A fetus that is larger than expected, often estimated at 4,000–4,500 grams or more at birth depending on the clinical definition. Diabetes and previous large babies increase risk; complications include difficult birth, injury, and cesarean delivery.

\medterm{Fetal Malpresentation} A fetal presentation other than the usual vertex presentation, including breech, transverse, or oblique lie, which may affect labor and delivery planning.

\medterm{Fetal Monitoring} Assessment of fetal heart rate and, sometimes, uterine contractions to evaluate fetal wellbeing during pregnancy or labor.

\medterm{Fetal Monitoring, Electronic} Continuous or intermittent recording of fetal heart rate and uterine activity using external sensors or an internal electrode when indicated; patterns are interpreted with gestational age and labor context.

\medterm{Fetal Presentation} The part of the baby positioned to enter the pregnant person’s pelvis first. The usual presentation is head first; breech, shoulder, face, brow, and mixed presentations can affect labor and delivery planning.

\medterm{Fetal Pyelectasis} Dilation of the fetal renal pelvis seen on prenatal imaging; it may be transient or associated with urinary tract obstruction or aneuploidy and requires context-specific follow-up.

\medterm{Fetal Scalp Electrode} A small electrode placed on the fetus’s scalp through the cervix to record the heart rate directly during labor. It may give a clearer signal when external monitoring is inadequate but can cause infection, bleeding, or scalp injury.

\medterm{Fetal Scalp PH Testing} A test that analyzes a small blood sample from the baby's scalp during labor to assess whether reduced oxygen may be causing acid buildup. It may clarify an abnormal fetal heart-rate pattern, but use is limited and depends on local expertise and clinical circumstances.

\medterm{Fetal Scalp Stimulation} A test during labor in which the baby’s scalp is gently touched during a vaginal examination. A brief rise in the baby’s heart rate may be reassuring, but the result must be interpreted with the full heart-rate recording and clinical situation.

\medterm{Fetal Stage} The fetal stage begins after the embryonic period and continues until birth, characterized mainly by growth and maturation of established organs.

\medterm{Fetal Tachycardia} A sustained fetal heart-rate baseline above 160 beats per minute for at least 10 minutes. Causes may include maternal fever, infection, medicines, fetal hypoxia, or fetal arrhythmia.

\medterm{Fetal Tissue} Fetal tissue is tissue from a developing fetus and may be relevant to pregnancy care, pathology, research, transplantation, or genetic testing under applicable ethical and legal rules.

\medterm{Fetal Ultrasound} Imaging during pregnancy using sound waves to assess the fetus, placenta, fluid, growth, movement, and anatomy. It can estimate gestational age, identify some structural problems, assess growth or blood flow, and guide selected procedures.

\medterm{Fetal Viability} The possibility that a fetus could survive outside the uterus with intensive medical care. It depends on pregnancy age, fetal condition, treatment available, and other factors and is not defined by one fixed week for every pregnancy.

\medterm{Fetal Weight} An estimate of a fetus’s mass, usually calculated from ultrasound measurements. It helps assess growth and identify possible growth restriction or unusually large size, but it is not exact.

\medterm{Fetal-Maternal Erythrocyte Distribution Blood Test} A test that estimates the amount of fetal red blood cells in maternal circulation, helping guide Rh immune globulin dosing after fetomaternal hemorrhage.

\medterm{Fetishism} A recurrent sexual interest in a nonliving object or a specific non-genital body part; it is a disorder only when it causes clinically significant distress or impairment or involves harm or lack of consent.

\medterm{Fetishistic Disorder} A mental-health disorder involving repeated, intense sexual arousal from a nonliving object or a specific non-genital body part. It is diagnosed only when it causes substantial distress or impairment or involves harm or lack of consent.

\medterm{Fetomaternal Hemorrhage} Passage of fetal blood into the pregnant person’s bloodstream through the placenta. Small amounts can occur normally, but heavy bleeding across the placenta after trauma, placental separation, or a procedure can cause fetal anemia and affect preventive treatment.

\medterm{Fetomaternal Transfusion} Movement of fetal blood cells across the placenta into the pregnant person's bloodstream. Small amounts are common, but heavy transfer after trauma, placental problems, or procedures can cause fetal anemia and may require urgent testing and treatment.

\medterm{Fetoscope} An instrument used to hear or assess the fetal heartbeat through the pregnant person's abdominal wall, especially later in pregnancy.

\medterm{Fetus} A developing human from the start of the ninth week after fertilization until birth. During this period, established organs grow and mature; medical descriptions may instead use pregnancy weeks counted from the last menstrual period.

\medterm{FEV1/FVC Ratio} The amount of air forcefully exhaled in the first second divided by the total amount exhaled. A lower-than-expected ratio suggests narrowed airways, but interpretation uses age, sex, height, reference values, and bronchodilator response.

\medterm{Fever} A rise in body temperature caused by the body’s response to infection, inflammation, certain medicines, heat illness, cancer, or autoimmune disease.

\medterm{Fever And Headache} Fever with headache may occur with viral illness, influenza, sinus disease, dehydration, meningitis, encephalitis, or other infection. Severe sudden headache, neck stiffness, confusion, rash, seizure, weakness, or reduced consciousness requires emergency assessment.

\medterm{Fever Blister} A small, painful blister on or around the lips caused by herpes simplex virus, usually type 1. The virus remains in nerve tissue and may reactivate with illness, stress, or sunlight. Blisters usually heal in about one to two weeks and can spread by close contact.

\synonyms
Cold sore; herpes labialis

\medterm{Fever Of Unknown Origin} Persistent or recurrent fever whose cause remains unexplained after an appropriate initial evaluation. Causes may include infection, inflammation, cancer, or medicines and require systematic reassessment.
\medterm{Fever, Acute, In Adults} A fever in an adult with recent onset, commonly caused by infection, inflammation, medication, heat illness, or other conditions. Severity, vital signs, immune status, source, and associated symptoms guide evaluation; confusion, breathing difficulty, stiff neck, shock, or severe localized pain requires urgent care.

\medterm{Fever, Chronic} A fever that continues or repeatedly returns over a long period. Causes include infection, inflammatory disease, cancer, and medicines; persistent unexplained fever needs medical evaluation based on its pattern and accompanying symptoms.

\medterm{Fever, Valley} Coccidioidomycosis: a fungal infection caused by Coccidioides species, endemic to arid regions of the southwestern US and parts of Mexico and Central/South America. Inhalation of spores causes respiratory illness ranging from asymptomatic to severe pneumonia. Dissemination can affect skin, bones, and meninges.

\synonyms
Valley fever; coccidioidomycosis

\medterm{Fevers, Viral Hemorrhagic} A group of serious viral illnesses that can damage blood vessels and affect several organs. Some cause bleeding, shock, or organ failure. They spread in different ways, including through mosquitoes, ticks, rodents, infected animals, or people; urgent isolation and specialist care may be needed.

\synonyms
Viral hemorrhagic fevers

\medterm{Fexofenadine} An antihistamine used to relieve sneezing, itching, runny nose, and watery eyes from allergies and itching from hives. It usually causes less drowsiness than older antihistamines, but headache or nausea can occur.

\medterm{Fiber} Dietary fiber is plant material that the small intestine cannot digest. It adds bulk to stool, supports regular bowel movements, and some types help lower blood cholesterol and slow rises in blood glucose.

\medterm{Fiber Optic} A technology that sends light through thin glass or plastic fibers. Flexible fiber-optic instruments are used in endoscopy to view body passages and guide some minimally invasive procedures.

\medterm{Fibrates} A group of medicines that mainly lower high triglyceride levels by changing how the liver and tissues handle fats. They can cause indigestion, gallstones, liver problems, or muscle injury, especially with some other lipid medicines.

\medterm{Fibre} A threadlike structure in the body or another material. In nutrition, fibre is plant material that is not digested in the small intestine and helps regulate bowel movements.

\medterm{Fibrillation} Rapid, disorganized activity of muscle fibers. In the heart, atrial fibrillation causes an irregular rhythm and can lead to blood clots; ventricular fibrillation stops effective pumping.

\medterm{Fibrin} An insoluble, threadlike protein that forms a mesh at an injured blood vessel. The mesh traps blood cells and strengthens the clot that helps stop bleeding.

\medterm{Fibrin Degradation Products Blood Test} A blood test that measures fragments released when blood clots are broken down. High levels may occur with widespread clotting, severe infection, liver disease, recent surgery, or pregnancy, so results require clinical interpretation.

\medterm{Fibrin Sealant} A surgical adhesive containing clotting proteins that form a fibrin clot when applied to tissue. It can help control bleeding or seal selected tissues, but it does not replace repair of major blood-vessel injuries.

\medterm{Fibrinogen} A blood protein made mainly by the liver that is changed into fibrin during clot formation. Low levels can cause bleeding; levels may rise with inflammation and fall in severe liver disease or widespread clotting.

\medterm{Fibrinogen Assay} A laboratory test that measures the amount or clotting activity of fibrinogen in blood. It helps investigate unexplained bleeding, abnormal clotting, liver disease, or suspected widespread clotting disorders.

\medterm{Fibrinogen Blood Test} A blood test measuring fibrinogen, a liver-made protein needed to form blood clots. Low levels may contribute to bleeding, while high levels commonly accompany inflammation, infection, pregnancy, or cardiovascular risk.

\medterm{Fibrinogen Decreased} A lower-than-normal blood fibrinogen level, which can weaken clot formation and cause bleeding. Causes include inherited deficiency, severe liver disease, major blood loss or transfusion, and widespread activation of clotting.

\medterm{Fibrinoid Necrosis} Severe damage to a blood-vessel wall in which blood proteins and immune material collect within the wall. It can occur in some inflammatory vessel diseases and is identified by examining affected tissue.

\medterm{Fibrinolysis} The natural process that dissolves a blood clot after healing begins. An enzyme called plasmin breaks down fibrin; excessive breakdown can cause bleeding, while insufficient breakdown can allow harmful clots to persist.

\medterm{Fibrinolysis - Primary Or Secondary} Primary fibrinolysis is early clot breakdown without preceding widespread clot formation. Secondary fibrinolysis follows clotting activation, as the body breaks down fibrin formed during a clotting disorder or injury.

\medterm{Fibrinolytic Agent} A medicine that dissolves a blood clot by helping the body convert plasminogen into plasmin, which breaks down fibrin. It may treat selected emergencies but can cause serious bleeding, including bleeding in the brain.

\medterm{Fibrinolytic Therapy} Treatment that uses clot-dissolving medicine to restore blood flow in selected emergencies, such as certain strokes, heart attacks, or pulmonary embolisms. The main risk is serious bleeding, so eligibility is assessed urgently.

\synonyms
Thrombolytic therapy

\medterm{Fibrinopeptide A Blood Test} A blood test that measures a protein fragment released when the body begins converting fibrinogen into fibrin. It can indicate clotting activity but is not specific enough to diagnose a condition by itself.

\medterm{Fibrinous Pericarditis} Inflammation of the sac surrounding the heart, with a sticky protein called fibrin coating its surfaces. It may cause sharp chest pain, a scratchy heart sound, and fluid around the heart.

\medterm{Fibroadenoma} A common noncancerous breast lump made of glandular and connective tissue. It is usually firm, smooth, movable, and painless, especially in adolescents and young adults; imaging helps distinguish it from other breast masses.

\medterm{Fibroadenoma Of The Breast} A noncancerous breast lump containing glandular and connective tissue. It is commonly smooth, firm, movable, and painless, but any new or changing breast lump should be examined and appropriately imaged.

\medterm{Fibroblast} A cell that makes collagen and other supporting materials in connective tissue. Fibroblasts help maintain skin, tendons, and organs and become active during wound healing and scar formation.

\medterm{Fibrocartilage} A tough type of cartilage that combines strength with limited flexibility. It helps absorb pressure and resist pulling in the discs between spinal bones, the pubic joint, and the knee menisci.

\medterm{Fibrocystic Breast Changes} Common noncancerous breast changes involving thickened tissue or fluid-filled cysts. Breast fullness, lumpiness, and tenderness often vary with hormonal changes during the menstrual cycle.

\medterm{Fibrocystic Breast Condition} A common noncancerous pattern of breast lumpiness, tenderness, or cysts that may change during the menstrual cycle. A persistent or changing lump, skin change, nipple inversion, or bloody discharge needs medical assessment.

\medterm{Fibrocystic Breast Disease} An older name for common noncancerous breast changes causing lumpiness, swelling, cysts, or tenderness. Symptoms may vary with the menstrual cycle, and the changes themselves do not increase breast-cancer risk.

\medterm{Fibrocystic Breasts} Common noncancerous changes in breast tissue that may cause lumpiness, thickening, cysts, swelling, tenderness, or discomfort. Symptoms may change during the menstrual cycle and usually improve after menopause. Fibrocystic breast changes do not increase the risk of breast cancer.

\medterm{Fibrodysplasia Ossificans Progressiva} A rare inherited disorder in which painful episodes gradually turn muscles, tendons, and other soft tissues into bone. This progressive extra bone can severely restrict movement and complicate medical procedures.

\synonyms
FOP; Stone Man Syndrome

\medterm{Fibroepithelial Polyp Of The Vulva} A usually noncancerous, soft tissue growth on the vulva or nearby vaginal tissue. It may enlarge during pregnancy or the reproductive years and can cause irritation, bleeding, or discomfort if large.

\medterm{Fibrofolliculoma} A small, usually noncancerous skin growth arising around a hair follicle, often on the face, neck, or upper trunk. Multiple lesions may occur with an inherited disorder that also raises kidney-tumor risk.

\medterm{Fibroid} A noncancerous growth of smooth muscle, most often in the uterus. It may cause heavy menstrual bleeding, pelvic pressure, pain, frequent urination, fertility problems, or no symptoms.

\synonyms
Uterine leiomyoma

\medterm{Fibroid Tumor} A noncancerous growth of the uterus made from smooth muscle and connective tissue. It may cause heavy or painful periods, pelvic pressure, frequent urination, infertility, or no symptoms.

\medterm{Fibroids} Noncancerous growths in the muscle of the uterus. They may cause heavy or painful menstrual bleeding, pelvic pressure, frequent urination, constipation, fertility problems, or no symptoms.

\medterm{Fibroids, Cysts, And Polyps} Three usually noncancerous types of growth: fibroids arise from uterine muscle, cysts are fluid-filled sacs often found in the ovaries, and polyps are tissue projections from the uterine lining or cervix.

\medterm{Fibroids, Uterine} Noncancerous smooth-muscle growths in the uterus. They may cause heavy bleeding, pelvic pain, pressure, frequent urination, or fertility problems; management depends on symptoms, size, location, age, and pregnancy plans.

\synonyms
Uterine fibroids; leiomyomas

\medterm{Fibrolipoma} A usually noncancerous, slow-growing lump made of mature fat and fibrous tissue. It is often painless and lies beneath the skin, but a new or enlarging lump should be assessed.

\medterm{Fibroma} A usually noncancerous tumor made of fibrous connective tissue. Cardiac fibromas most often occur in a child’s heart muscle or dividing wall and may obstruct blood flow or disturb the heartbeat.

\medterm{Fibroma Of Tendon Sheath} A usually noncancerous, slow-growing lump attached to the covering of a tendon, most often in a finger, hand, or wrist. It is firm and may recur after removal.

\medterm{Fibromatoses} A group of connective-tissue growth disorders that do not spread to distant organs but may grow into nearby structures and return after treatment. Effects depend on the affected body site.

\medterm{Fibromatosis} An overgrowth of connective-tissue cells that does not spread to distant organs but can invade nearby tissue and recur. Desmoid-type disease may cause pain, a firm mass, or impaired organ function.

\medterm{Fibromuscular Dysplasia} An abnormal structure of medium-sized artery walls, most often in the kidney or neck arteries. It can cause high blood pressure, headache, or stroke, or produce an artery bulge or tear.

\medterm{Fibromyalgi} Alternate index label; misspelling of \textit{Fibromyalgia}.

\medterm{Fibromyalgia} A long-term condition involving increased sensitivity to pain, with widespread aching, fatigue, unrefreshing sleep, and problems with concentration or memory. Diagnosis is based on symptoms and examination; no single laboratory test confirms it.


\medterm{Fibromyalgia Syndrome} A long-term pain condition causing widespread body pain, fatigue, poor-quality sleep, and difficulty thinking or concentrating. Doctors diagnose it from the symptom pattern and examination, while treatment combines activity, sleep support, and selected medicines.

\synonyms
Fibromyalgia

\medterm{Fibromyxosarcoma} A rare cancer of soft tissue containing fibrous and jelly-like material. It may present as a growing lump or pain; diagnosis requires tissue examination, and treatment depends on its site and spread.

\medterm{Fibronectin} A protein found in blood and connective tissue that helps cells attach, move, and organize during tissue repair. It supports wound healing and the structure of tissues around cells.

\medterm{Fibroproliferation} Excess growth of connective-tissue cells with increased production of scar tissue. It can thicken or stiffen organs, including the lungs, and may develop after injury, inflammation, or surgery.

\medterm{Fibrosarcoma} A rare cancer arising from cells that make fibrous connective tissue. It usually develops in deep soft tissue or bone and may cause a growing mass, pain, or reduced movement.

\medterm{FibroScan} A painless ultrasound-based test that estimates liver stiffness. Greater stiffness can suggest scarring, but inflammation, congestion, obesity, or fluid in the abdomen can affect accuracy and results may need other tests.

\medterm{Fibrosing Alopecia} Hair loss caused by inflammation and permanent scarring of hair follicles. The affected skin may look smooth or shiny, and regrowth is unlikely once follicles are destroyed; early treatment may limit progression.

\medterm{Fibrosis} Excess formation of scar tissue after ongoing injury or inflammation. Fibrosis can make an organ stiff and reduce its function, such as causing breathlessness in the lungs or impaired filtration in the kidneys.

\medterm{Fibrosis, Interstitial Pulmonary} Scarring in the supporting tissue between the lung air sacs, which can reduce oxygen transfer. It may cause progressive breathlessness and dry cough; causes include autoimmune disease, medicines, radiation, or no identified cause.

\synonyms
Interstitial lung disease

\medterm{Fibrosis, Pulmonary} Scarring of lung tissue that makes the lungs stiff and reduces oxygen exchange. It often causes gradually worsening breathlessness and dry cough; examination, lung tests, and scans help identify the cause and extent.

\synonyms
Pulmonary fibrosis

\medterm{Fibrositis} An outdated term formerly used for muscle and soft-tissue pain. It is generally avoided because true inflammation is usually not shown and a more specific diagnosis, such as fibromyalgia, may be appropriate.

\medterm{Fibrous Dysplasia} A noncancerous bone disorder in which normal bone is replaced by weaker abnormal tissue. It may affect one or several bones and cause pain, deformity, fractures, or pressure on nearby nerves.

\medterm{Fibrous Plaque} A thickened area in an artery wall containing scar-like tissue, cholesterol, and other material. It can narrow blood flow or rupture, allowing a blood clot to form and block the artery.

\synonyms
Atherosclerotic plaque

\medterm{Fibula} The slender bone on the outer side of the lower leg. It provides muscle attachment and forms the outer ankle. The nerve near its upper end can be injured by pressure or fracture.

\medterm{Fick Principle} A method for estimating blood flow from how much oxygen the body uses and the difference in oxygen content between blood entering and leaving the lungs. It can estimate cardiac output during specialized heart testing.

\synonyms
Fick Method; Fick Cardiac Output Determination

\medterm{Fiducial Markers} Small visible markers placed on or inside the body to identify a precise location during imaging, radiation treatment, or surgery. They help clinicians align equipment and target the intended tissue accurately.

\medterm{Field Fever} An older name for leptospirosis, a bacterial infection acquired through water or soil contaminated with urine from infected animals. It may cause fever, headache, muscle pain, jaundice, kidney injury, or meningitis.

\medterm{Field Of View} The area included in a medical image or examination. A wider field shows more of the body, while a narrower field provides greater detail about a smaller area.

\medterm{Fifth Disease} A usually mild infection caused by parvovirus B19. Children may develop fever, cold-like symptoms, bright-red cheeks, and a lacy rash. It can cause serious fetal anemia during pregnancy and prolonged anemia in some vulnerable people.

\synonyms
Erythema Infectiosum (Fifth Disease)

\medterm{FIGO Leiomyoma Classification} A system that describes where a uterine fibroid lies in relation to the uterine cavity and outer surface. The location category helps clinicians predict symptoms and choose treatment.

\medterm{Filamentous Fungus} A fungus that grows as branching threads called hyphae, which form a network called mycelium. Some filamentous fungi cause skin, nail, lung, or invasive infections, especially in people with weakened immunity.
\medterm{Filarial Elephantiases} Severe enlargement and thickening of body parts caused by long-term lymphatic obstruction from filarial worm infection. It most often affects the legs, arms, breasts, or genitals and develops gradually.

\medterm{Filarial Elephantiasis} Long-lasting swelling and thickening of skin and underlying tissues caused by lymphatic blockage from filarial worms. It commonly affects the legs or genitals and may require antiparasitic treatment, skin care, and swelling management.

\medterm{Filariasis} Infection with threadlike parasitic worms spread by mosquitoes, flies, or other insects. Depending on the species, worms may affect lymph vessels, skin, or eyes, causing swelling, inflammation, itching, or vision problems.

\medterm{Filler} A material injected or placed in tissue to restore volume, smooth wrinkles, change contours, or replace lost structure. Risks include pain, infection, lumps, and accidental blockage of a blood vessel.

\medterm{Filler Excipient} An inactive ingredient added to a medicine to provide bulk, improve stability, or help form a tablet, capsule, or liquid. It may rarely cause allergy or intolerance.

\medterm{Filopodium} A thin, finger-like projection from a cell surface. It senses nearby surroundings and helps cells attach, move, communicate, or form connections, including connections made by developing nerve cells.

\medterm{Filoviridae} A family of viruses that includes Ebola and Marburg viruses. They can cause severe fever, bleeding, organ failure, and death, and are mainly spread through direct contact with blood or other infected body fluids.

\medterm{Filoviruses} Viruses such as Ebola and Marburg that can cause severe hemorrhagic fever. Symptoms may include fever, headache, muscle pain, vomiting, diarrhea, bleeding, and organ failure; spread occurs through infected body fluids.

\medterm{Filter (Vascular)} A small device placed in a large vein, usually the inferior vena cava, to catch clots traveling toward the lungs. It is used selectively when anticoagulant treatment is unsuitable or ineffective.

\synonyms
Vena cava filter; inferior vena cava filter

\medterm{Filtered Back Projection} A computer method for creating an image from scan measurements, especially in computed tomography. It is fast, but images may contain more noise or artifacts than those made with newer reconstruction methods.

\medterm{Filtering Procedure} Eye surgery that creates a new drainage route for fluid to lower pressure inside the eye. It may treat glaucoma when medicines or other procedures do not control pressure adequately.

\medterm{Filtration} Separation of particles or dissolved substances from a fluid by passing it through a barrier. In the kidneys, tiny filters remove water and small waste products from blood to begin urine formation.

\medterm{Filtration Fraction} The percentage of plasma flowing through the kidneys that is filtered to begin urine formation. It helps describe kidney blood flow and filtering function and is calculated from kidney filtration compared with plasma flow.

\medterm{Filum Terminale} A thin strand of connective tissue that anchors the lower end of the spinal cord to the tailbone. Abnormal attachment can restrict cord movement and cause symptoms of tethered spinal cord.

\medterm{Fimbria} A fringe-like structure. In the fallopian tube, the fimbria are finger-like projections near the ovary that help guide a released egg into the tube; the term is also used for bacterial attachment structures.

\medterm{Fimbriae} Finger-like projections at the end of each fallopian tube that help collect an egg released from the ovary and guide it into the tube. Their movement and tiny surface hairs help direct the egg toward the uterus.

\medterm{Fimbriae Proteins} Bacterial surface proteins that form hair-like projections used to attach bacteria to human cells, medical devices, or other surfaces. Attachment can help bacteria form communities and begin infection.

\medterm{Fine Crackles} Brief, high-pitched, popping sounds heard mainly during breathing in. They may occur when small airways or air sacs open in conditions such as lung scarring, fluid congestion, or other lung disease.

\medterm{Fine Motor Control} The coordinated use of small muscles for accurate movements, especially of the hands and fingers. It supports tasks such as writing, buttoning clothes, using utensils, and handling small objects.

\medterm{Fine Motor Skills} Abilities that use small muscles for controlled movements, especially in the hands and fingers. Children develop skills such as grasping, feeding themselves, drawing, dressing, and writing at different individual rates.

\medterm{Fine Needle Aspiration} A procedure using a thin needle to remove cells or fluid from a lump or organ for laboratory examination. It usually causes little tissue injury and may be guided by touch or ultrasound.

\medterm{Fine Needle Aspiration Of The Thyroid} A procedure using a thin needle, often guided by ultrasound, to collect thyroid cells from a nodule. Laboratory examination helps determine whether the nodule is benign, malignant, or needs further testing.

\medterm{Finegoldia} A group of bacteria that grow without oxygen and normally live on skin or mucous membranes. They can occasionally cause wound, soft-tissue, bone, joint, or bloodstream infections.

\medterm{Finegoldia Magna} A bacterium normally found on skin and mucous membranes that grows without oxygen. It can cause infections after wounds or surgery, including abscesses and deeper soft-tissue infections.

\medterm{Finger} One of the hand’s digits, made of bones, joints, muscles, tendons, nerves, blood vessels, and skin. Fingers support grasping, touch, balance, and precise movements.

\medterm{Finger Dislocation} An injury in which the bones of a finger joint are forced out of their normal alignment. Pain, swelling, deformity, and limited movement are common; fractures, tendon damage, and impaired circulation must also be checked.

\medterm{Finger Flexors} Muscles and tendons that bend the fingers at their joints. Injury can cause pain or inability to bend a finger and may require prompt assessment, especially after a cut or forceful injury.

\medterm{Finger Injury} Damage to a finger’s bone, joint, tendon, ligament, nerve, blood vessel, nail, or skin, usually from trauma. Assessment checks alignment, movement, sensation, blood flow, and possible hidden fractures.

\medterm{Finger Joint} A connection between the bones of a finger that allows bending, straightening, and limited side-to-side movement. Finger joints can be affected by injury, arthritis, inflammation, infection, or dislocation.

\medterm{Finger Pain} Discomfort in one or more fingers caused by injury, arthritis, tendon problems, nerve compression, infection, inflammation, or reduced blood flow. Sudden severe pain, deformity, numbness, or discoloration needs prompt care.

\medterm{Finger Stick} A small fingertip puncture used to collect a drop of capillary blood. The sample can be tested immediately, commonly for blood glucose or other point-of-care measurements.

\medterm{Fingernail} The hard protective plate covering the end of a finger, made mainly of keratin and growing from the nail root. Changes in color, shape, or growth can sometimes signal injury or disease.

\medterm{Fingers That Change Color} Fingers that become white, blue, purple, or red because of changes in blood flow, temperature, or oxygen. Recurrent attacks may be Raynaud phenomenon; persistent changes, pain, numbness, or sores need medical assessment.

\medterm{Finkelstein Test} A wrist examination that gently stresses the tendons on the thumb side. Pain there may support de Quervain tenosynovitis, but the result is not specific and must be interpreted with symptoms and examination.

\medterm{Fire Ants} Ants whose stings inject venom and cause immediate burning pain, redness, and itchy pus-filled bumps. Multiple stings can cause severe swelling, vomiting, breathing difficulty, or a life-threatening allergic reaction.


\medterm{Firesetting Behavior} Deliberately starting fires. Repeated fire-setting may be linked to an impulse-control disorder, mental illness, substance use, developmental problems, or social stress and requires assessment of safety and risk to others.

\medterm{First Aid} Immediate basic care for injury or sudden illness before professional treatment arrives. It includes ensuring safety, calling emergency services, checking breathing, starting resuscitation when needed, and controlling severe bleeding.

\medterm{First Aid Kit} A group of supplies for treating minor injuries and giving temporary care during emergencies. Useful items include gloves, dressings, bandages, antiseptic, scissors, and emergency contact information.

\medterm{First Stage Of Labor} The time from regular labor contractions that cause cervical change until the cervix is fully open. It includes early and active phases, with progress varying according to the person, fetus, and clinical circumstances.

\medterm{First Trimester} The first twelve weeks of pregnancy, when the embryo’s organs begin forming and the placenta develops. Nausea, fatigue, and breast tenderness are common; prenatal care includes dating and early health assessment.

\medterm{First-Degree AV Block} A mild delay in the electrical signal traveling from the heart’s upper chambers to its lower chambers. Every signal still reaches the lower chambers; it is often symptomless but may reflect medicines or heart disease.

\medterm{First-Degree Burn} A superficial burn affecting only the outer skin layer. It causes redness, pain, and dryness without blisters, as with mild sunburn, and usually heals within several days without scarring.

\medterm{First-Degree Heart Block} A delay in the electrical signal between the heart’s upper and lower chambers, shown on an electrocardiogram by a prolonged but constant interval. It is often harmless, though medicines or heart disease may contribute.

\medterm{First-Episode Psychosis} The first recognized period of losing contact with reality, such as hearing or seeing things that are not present or holding fixed false beliefs. Urgent assessment looks for medical, medication, substance, and psychiatric causes.

\medterm{First-Pass Effect} The reduction of an orally swallowed medicine by the intestine and liver before it reaches the bloodstream. This can reduce the amount available to act and may make another route necessary.

\medterm{First-Pass Metabolism} Breakdown of a swallowed medicine by the intestinal wall and liver before it enters general circulation. Extensive breakdown lowers the active amount reaching the body and can influence dose or route of administration.

\medterm{Fish Oil} Oil containing omega-3 fatty acids, usually obtained from oily fish. Supplements may lower triglycerides at prescribed doses, but benefits vary and side effects or medicine interactions should be reviewed with a clinician.

\medterm{Fish Poisoning} Illness caused by toxins or germs in contaminated or improperly stored fish. Symptoms may include vomiting, diarrhea, flushing, tingling, weakness, or breathing difficulty; severe symptoms require urgent medical care.

\medterm{Fish Tapeworm Infection} An intestinal infection acquired by eating raw or undercooked fish containing tapeworm larvae. It may cause abdominal discomfort, diarrhea, weight loss, or vitamin B12 deficiency and is usually treatable with antiparasitic medicine.

\medterm{Fish Venom} Poison produced by certain fish, usually delivered through spines or sharp body parts. It can cause immediate severe pain, swelling, numbness, tissue injury, or rarely whole-body poisoning.

\medterm{Fish-Handler'S Disease} A skin infection caused by Erysipelothrix bacteria after handling infected fish, shellfish, meat, or animals. It usually causes a slowly expanding, painful purple-red lesion on the hand or fingers.

\medterm{Fisher'S Syndrome} A rare nerve disorder, usually after an infection, causing eye-muscle weakness, poor coordination, and reduced or absent reflexes. It is a variant of Guillain-Barré syndrome and needs prompt medical assessment.

\medterm{Fishhook Removal} Removal of a fishhook from skin or deeper tissue. The method depends on the hook’s location and depth; clinicians assess nerves, blood vessels, tendons, joints, tetanus protection, and infection risk.

\medterm{Fissure} A narrow crack, groove, or split in a body structure. An anal fissure is a painful tear in the anal lining, often causing sharp pain and bright-red bleeding during bowel movements.

\medterm{Fissured Tongue} A usually harmless condition in which grooves develop on the tongue’s surface. Food may collect in the grooves and cause irritation or odor; gentle tongue cleaning and routine oral care can help.

\medterm{Fistula} An abnormal passage connecting two organs or a body organ to the skin. It may develop because of infection, inflammation, surgery, injury, cancer, or abnormal development and can cause persistent drainage.

\medterm{Fistula (Pathology)} Alternate index label; see \textit{Fistula}.

\medterm{Fistula, Anal} An abnormal tunnel between the anal canal or rectum and nearby skin, often left after an abscess. It may cause persistent drainage, pain, skin irritation, or repeated infections and often needs surgical treatment.

\medterm{Fistula, Arteriovenous} Alternate word order; see \textit{Arteriovenous Fistula}.

\medterm{Fistula-In-Ano} An abnormal tunnel between the anal canal and skin near the anus. It commonly follows an anal abscess and may occur with Crohn disease, causing ongoing drainage, discomfort, or recurrent abscesses.

\medterm{Fit} A seizure or sudden episode of altered movement or awareness; this everyday term is also used to mean physically healthy. In medical writing, the specific symptom or diagnosis should be stated instead.

\medterm{Fits (Seizures)} Episodes caused by abnormal electrical activity in the brain. They may involve staring, unusual sensations, jerking, stiffening, loss of awareness, or unconsciousness; a first or prolonged seizure needs urgent assessment.

\synonyms
Seizures

\medterm{Fitz-Hugh-Curtis Syndrome} Inflammation and scarring around the liver caused by pelvic inflammatory disease, usually from gonorrhea or chlamydia. It can cause sharp right-upper-abdominal pain, fever, and pain that worsens with movement or breathing.

\medterm{Five-Year Survival Rate} The percentage of people alive five years after a diagnosis or treatment. It summarizes outcomes for groups, not an individual’s exact prognosis, which also depends on cancer type, stage, health, and treatment.

\medterm{Fixation} Holding a body part, tissue, or broken bone in the correct position while it heals. Methods include stitches, wires, screws, plates, rods, casts, splints, or external frames.

\medterm{Fixation, Internal} Surgery that stabilizes a broken or unstable bone with devices placed inside the body, such as plates, screws, pins, or rods. The aim is to restore alignment and allow healing and movement.

\medterm{Fixative} A chemical used to preserve cells or tissue for laboratory examination by preventing decay and maintaining structure. Common fixatives can irritate skin, eyes, and airways and must be handled safely.

\medterm{Fixed Joint} A joint that permits little or no movement because the bones are joined by bone or dense connective tissue. Some fixed joints are normal; others result from injury, disease, or surgical fusion.
\synonyms
Synarthrosis

\medterm{Fixed Macrophage} An immune cell that stays within a particular tissue instead of circulating in blood. It removes dead cells and microbes and helps coordinate inflammation and repair.

\medterm{Flaccid} Soft and lacking normal firmness or muscle tone. It may describe a limb, muscle, bladder, or penis; sudden flaccid weakness, especially on one side, can signal a neurologic emergency.

\medterm{Flagellin} The main protein in a bacterial flagellum, the whip-like structure that helps some bacteria move. Human immune cells can recognize flagellin and respond with inflammation.

\medterm{Flagellum} A long, whip-like structure that helps a cell or microorganism move. Human sperm and some bacteria have flagella, which differ in structure and function.
\medterm{Flail Chest} A serious chest injury in which several ribs break in more than one place, leaving part of the chest wall unstable. Severe pain and difficult breathing can result, often with lung bruising and need for urgent care.

\medterm{Flakka} A street name for illicit synthetic stimulant drugs, especially alpha-PVP. Overdose can cause extreme agitation, hallucinations, seizures, dangerously high temperature, muscle breakdown, irregular heartbeat, or death and requires emergency treatment.

\medterm{Flanged Needle} A needle with a widened rim or attached flange that helps stabilize it, limit how far it enters, or connect it to another device. The flange may reduce movement and tissue injury during sampling, injection, drainage, or infusion.

\medterm{Flank} The side of the body between the lower ribs and the hip, including part of the abdomen and back. Pain in this area may arise from the kidney, muscles, spine, bowel, or nearby organs.

\medterm{Flank Pain} Pain between the lower ribs and hip, often caused by kidney stones, kidney infection, muscle strain, spine problems, or nearby abdominal disease. Fever, vomiting, blood in urine, injury, or severe pain needs prompt evaluation.

\medterm{Flap Operation} Surgery that moves a piece of skin, tissue, or muscle, while preserving its blood supply, to cover a wound or reconstruct a damaged body part. Risks include infection, bleeding, and loss of the flap.

\medterm{Flare} A temporary increase in symptoms or inflammation after a disease has been stable or less active. Flares may be triggered by infection, stress, injury, or treatment changes; severe or unusual symptoms need assessment.

\synonyms
Exacerbation

\medterm{Flare-Up} A sudden or gradual worsening of a long-term disease after a period of relative stability. The cause may be infection, stress, missed treatment, or disease progression and should be assessed if severe or persistent.

\medterm{Flashes} Brief flickers or streaks of light seen without an outside light source. They may result from changes in the eye’s gel, migraine, or inflammation; sudden flashes with floaters or missing vision require urgent eye assessment.

\medterm{Flashing} Seeing brief flashes or streaks of light without an external source. New flashing, especially with floaters, a curtain-like shadow, eye pain, or reduced vision, can indicate a retinal emergency.

\medterm{Flat Affect} Very little outward expression of emotion, such as reduced facial movement, gestures, or voice variation. It is a clinical sign, not a diagnosis, and may occur with psychiatric, neurologic, or medication-related conditions.

\medterm{Flat Bones} Broad, relatively thin bones such as the skull, ribs, breastbone, and shoulder blades. They protect organs, provide surfaces for muscle attachment, and contain marrow that produces blood cells.

\medterm{Flat Feet} A lowered or absent arch along the inner side of the foot, causing more of the sole to touch the ground. It is often painless, but rigid or painful flat feet may reflect tendon, joint, or bone problems.

\medterm{Flatfoot} A foot with a reduced or collapsed inner arch. It may be flexible and painless or rigid and painful; causes include normal development, loose ligaments, tendon problems, injury, arthritis, or bone shape.

\medterm{Flatulence} Passage of gas through the anus, produced by swallowed air and bacterial digestion of food in the intestine. It may increase with diet, constipation, swallowing air, or digestive disorders; persistent troublesome symptoms need assessment.

\medterm{Flatus} Gas released through the anus after air is swallowed or intestinal bacteria break down food. Passing gas is normal; inability to pass gas with increasing abdominal swelling, vomiting, or pain may indicate bowel obstruction.

\synonyms
Intestinal gas

\medterm{Flavin Adenine Dinucleotide} A helper molecule made from vitamin B2 that allows enzymes to transfer electrons during energy production and other chemical reactions. It is commonly abbreviated FAD.

\synonyms
Flavin Adenine Dinucleotide

\medterm{Flaviviridae} A family of single-stranded RNA viruses that includes dengue, yellow fever, Zika, hepatitis C, and West Nile viruses. Different members spread through insects, blood, or other routes and cause different illnesses.

\medterm{Flavivirus} A group of RNA viruses that includes dengue, yellow fever, West Nile, Japanese encephalitis, and Zika viruses. Many spread through mosquitoes or ticks and can cause fever, inflammation, or nervous-system disease.

\medterm{Flaviviruses} Viruses in the flavivirus group, including dengue, Zika, yellow fever, and West Nile viruses. They are often spread by mosquitoes or ticks and may cause fever, bleeding, liver disease, birth defects, or brain inflammation.

\medterm{Flavobacteriaceae} A family of mainly gram-negative bacteria found in water, soil, and other environments. Most do not infect people, but some can cause opportunistic infections, especially in hospitalized or immunocompromised patients.

\medterm{Flavobacterium} A group of gram-negative bacteria commonly found in water and soil. A few species can cause opportunistic infections in people or animals, particularly when immunity is weakened.

\medterm{Flavoxate} A medicine used to relieve bladder spasms, urgency, frequency, or painful urination. It treats symptoms rather than infection or blockage and can cause dry mouth, blurred vision, dizziness, or constipation.

\medterm{Flaxseed Extract} A preparation made from flaxseed that may contain omega-3 fat, lignans, and soluble fiber. Effects vary by product; it may interact with medicines affecting clotting, blood sugar, or blood pressure.

\medterm{Flea} A small wingless insect that feeds on blood. Flea bites can cause itchy skin reactions, and fleas can transmit infections between animals or, less commonly, to people.

\medterm{Flea Bites} Itchy red bumps caused by fleas feeding on the skin, often clustered around the ankles or lower legs. Scratching can cause infection; treating pets and the environment helps prevent new bites.

\medterm{Flea Bites In Humans} Small itchy bumps, often around the ankles or lower legs, caused by flea feeding. Allergic reactions may produce larger areas of swelling; control requires treating infested animals, bedding, rooms, and symptoms.

\synonyms
Bites Flea (Flea Bites In Humans)
Ctenocephalides Flea Bite (Flea Bites In Humans)

\medterm{Fleas} Small wingless insects that feed on blood and infest animals, homes, or clothing. They can cause itchy bites, allergic skin disease, anemia in heavily infested animals, and transmission of some infections.

\medterm{Flecainide} A prescription antiarrhythmic medicine that slows electrical signals in the heart to prevent selected abnormal rhythms. It is not suitable for everyone and can worsen dangerous rhythms, especially with structural heart disease.

\medterm{Flesh-Eating Bug} An informal name for necrotizing fasciitis, a rapidly spreading infection of deep skin and soft tissue. Severe pain out of proportion to visible changes, swelling, fever, blisters, or skin discoloration requires emergency care.

\medterm{Flex} To bend a joint so the angle between connected bones becomes smaller. For example, bending the elbow brings the forearm toward the upper arm; pain, swelling, stiffness, or weakness can limit this movement.

\synonyms
Flexion

\medterm{Flexibility} The range of movement available at a joint or group of joints. It depends on bones, muscles, tendons, ligaments, joint health, nervous-system control, age, training, and pain.

\medterm{Flexible Nasolaryngoscopy} An examination using a thin flexible camera passed through the nose to view the nasal passages, throat, and voice box. It helps investigate hoarseness, swallowing difficulty, blockage, lesions, and vocal-cord movement.

\medterm{Flexible Sigmoidoscopy} An examination using a flexible camera to view the rectum and lower colon. It can investigate bleeding, diarrhea, or bowel symptoms and monitor known disease; preparation and biopsy may be needed.

\medterm{Flexion} Movement that bends a joint and decreases the angle between connected bones, such as bending the elbow, knee, hip, or fingers.

\medterm{Flexor} A muscle or tendon that bends a joint or body part. Flexors in the arm and hand bend the fingers and wrist; injury can cause pain, weakness, or loss of movement.

\medterm{Flexor Digitorum Profundus} A deep forearm muscle whose tendons bend the end joints of the fingers and also assist wrist movement. Its inner part is supplied by the ulnar nerve and its outer part by the median nerve.

\medterm{Flexor Digitorum Superficialis} A forearm muscle whose tendons bend the middle joints of the fingers and assist wrist movement. Damage to the muscle, tendon, or median nerve can cause pain and difficulty bending the fingers.

\medterm{Flexural} Relating to a skin fold or the inner surface of a bent joint, such as the elbow crease, back of the knee, armpit, groin, or fold beneath the breast.

\medterm{Flibanserin} A prescription medicine for persistent, distressing low sexual desire in some premenopausal women. It is taken regularly rather than as needed and can cause dizziness, sleepiness, low blood pressure, or fainting, especially with alcohol.

\medterm{Flight-Or-Fight Response} The body’s rapid response to perceived danger, releasing stress hormones that increase heart rate, breathing, alertness, and muscle readiness so a person can escape or confront a threat.
\synonyms
Fight-or-flight response; stress response

\medterm{Floater} A spot, thread, or cobweb-like shape that drifts across vision, usually caused by changes in the gel inside the eye. Sudden new floaters, flashes, or missing vision require urgent eye assessment.

\medterm{Floaters} Moving spots, threads, cobwebs, or rings in the field of vision caused by shadows from material in the eye’s vitreous gel. They are often age-related, but sudden onset can indicate bleeding or a retinal tear.

\medterm{Floating Kidney} Unusual downward movement of a kidney when a person stands, sometimes called nephroptosis. It is often symptomless but may cause position-related flank pain, urinary symptoms, or rarely impaired urine drainage.

\medterm{Floating Rib} One of the eleventh or twelfth ribs, attached to the spine but not connected to the breastbone. These ribs are less protected and can be injured by direct trauma.

\medterm{Floating Ribs} The eleventh and twelfth pairs of ribs, which attach to the spine but not the breastbone. They are relatively mobile and less protected; fractures can cause pain and may injure nearby organs.

\medterm{Flocculation} The joining of tiny particles suspended in a liquid into larger clumps. Laboratory tests may use this visible reaction to detect antibodies or other substances; it is also used in water treatment and pharmaceutical preparation.

\medterm{Floor Of The Mouth Cancer} Cancer arising in the tissues beneath the tongue, most often from surface cells. A persistent ulcer, lump, pain, bleeding, numbness, or difficulty swallowing or speaking needs prompt examination.

\synonyms
Cancer, Floor Of The Mouth

\medterm{Floppy Baby Syndrome} An informal description of an infant with unusually low muscle tone and a limp posture. Causes include brain, spinal cord, nerve, muscle, genetic, or metabolic disorders and require medical evaluation.

\medterm{Floppy Iris Syndrome} A cataract-surgery complication in which the iris becomes unusually loose, moves with fluid currents, and narrows the working space. It is associated with some prostate medicines and should be anticipated by the surgeon.

\medterm{Floppy Valve Syndrome} A condition in which the mitral valve’s leaflets bulge backward into the upper heart chamber when the heart contracts. It is often harmless but may cause palpitations, chest discomfort, or leakage through the valve.

\synonyms
Mitral valve prolapse

\medterm{Flora} The microorganisms living normally or temporarily in a body site or environment. Some protect against harmful organisms, but they can cause infection if they enter another body site or defenses are weakened.


\medterm{Flossing} Cleaning between teeth with dental floss or another interdental device to remove plaque and food that a toothbrush cannot reach. Gentle daily cleaning can help prevent gum disease and cavities.

\medterm{Flotation} A laboratory method that separates material by density. For example, a flotation test can concentrate parasite eggs from stool so they are easier to identify under a microscope.

\medterm{Flow Cytometry} A laboratory test that examines cells or particles one at a time as they pass through a laser beam. It can count cells and identify their features, helping diagnose blood cancers and immune disorders.

\medterm{Flow-Volume Loop} A graph showing how much air moves in and out of the lungs at different speeds during forceful breathing. Its shape can help identify airway narrowing, reduced lung size, or an upper-airway blockage.

\medterm{Flowmeter} A device that measures the rate at which a gas or liquid moves. In medicine, it may measure oxygen delivery, intravenous fluid flow, or breathing flow such as peak expiratory flow.

\medterm{Floxacillin} An antibacterial medicine used in some countries to treat infections caused by susceptible bacteria, especially staphylococci. It can cause allergic reactions, diarrhea, liver problems, or other antibiotic-related effects.

\medterm{Floxuridine} An anticancer medicine that interferes with DNA production, slowing or killing rapidly dividing cells. It is used only in selected treatment plans and may cause low blood counts, mouth sores, diarrhea, or other toxicity.

\medterm{Flu} An acute respiratory infection caused by influenza viruses. It commonly causes sudden fever, cough, sore throat, headache, muscle aches, and fatigue and can lead to pneumonia or worsen chronic diseases.


\synonyms
Influenza; influenza infection

\medterm{Flu, Swine} Influenza A infection mainly affecting pigs. Some strains can infect people, causing an illness like seasonal flu; severity varies, and human-to-human spread can occur when a strain adapts to spread efficiently.

\synonyms
H1N1 flu; swine influenza

\medterm{Flucloxacillin} An antibiotic used for selected bacterial infections, especially those caused by susceptible staphylococci. It can cause nausea, diarrhea, allergy, or liver injury and should be taken only as prescribed.

\medterm{Fluconazole} An antifungal medicine used for selected yeast and other fungal infections. It can interact with many medicines and may cause nausea, rash, or liver injury; treatment depends on the infection and patient.

\medterm{Flucytosine} An antifungal medicine converted inside susceptible fungi into substances that interfere with genetic-material production. It is used for selected serious infections, usually with another antifungal, and requires blood-count and liver monitoring.

\medterm{Fludarabine} A chemotherapy medicine that interferes with DNA production in abnormal blood cells. It treats selected blood cancers but can severely suppress immunity and blood counts, increasing the risk of serious infections.

\medterm{Fludarabine Phosphate} A form of fludarabine used in selected blood-cancer treatment plans. It interferes with DNA production and can cause prolonged low blood counts, immune suppression, infections, and neurologic or other serious effects.

\medterm{Fludeoxyglucose F-18} A radioactive form of glucose used in positron-emission tomography scans. Active tissues absorb more of it, helping doctors assess some cancers, brain disorders, and heart muscle; uptake is not specific for cancer.

\medterm{Fludrocortisone} A medicine that helps the body retain salt and water and raise blood pressure. It treats some adrenal-hormone deficiencies and selected orthostatic disorders; monitoring is needed for swelling, high blood pressure, and low potassium.

\medterm{Fluid Around The Heart} Extra fluid in the sac surrounding the heart. Small amounts may cause no symptoms, while rapid or large collections can restrict heart filling and cause breathlessness, low blood pressure, or life-threatening tamponade.

\synonyms
Pericardial effusion

\medterm{Fluid Balance} The relationship between water entering and leaving the body. The kidneys, hormones, thirst, breathing, sweat, and stool help regulate it; disruption can cause dehydration, swelling, or abnormal blood salts.

\medterm{Fluid Imbalance} Too much or too little body water or an abnormal level of salts such as sodium or potassium. Causes include vomiting, diarrhea, kidney or heart disease, hormone disorders, dehydration, or excess fluid treatment.

\medterm{Fluid Management} Planning and monitoring fluids, salts, and urine output to maintain circulation and organ function. Treatment may involve drinking, intravenous fluids, diuretics, or fluid restriction according to the underlying problem.

\medterm{Fluid Overload} Excess water and salt in the body, causing swelling, rapid weight gain, high blood pressure, or fluid in the lungs. It may result from heart, kidney, or liver disease or too much fluid treatment.

\medterm{Fluid Therapy} Treatment using fluids, often through a vein, to restore hydration, blood volume, or salt balance. The type and amount depend on the cause; excess fluid can worsen swelling, heart failure, or lung congestion.

\synonyms
Intravenous fluid therapy

\medterm{Fluke} A flat, leaf-shaped parasitic worm. Different species infect the liver, lungs, intestines, or blood vessels and are acquired through contaminated food, water, or contact with freshwater.

\medterm{Flumazenil} An antidote that can reverse sedation and slowed breathing caused by some benzodiazepine medicines. It can trigger seizures or sudden withdrawal and is not suitable for every overdose, especially mixed overdoses.

\medterm{Flunarizine} A medicine used in some countries to prevent migraine and treat selected balance disorders. It can cause sleepiness, weight gain, depression, or movement problems and should be prescribed with appropriate monitoring.

\medterm{Flunisolide} An inhaled or nasal corticosteroid used to reduce inflammation in asthma or allergic rhinitis. It may cause throat irritation, nosebleeds, or oral thrush; regular use does not relieve a sudden severe asthma attack.

\medterm{Fluocinolone} A corticosteroid applied to skin, scalp, or other affected areas to reduce inflammation and itching. Prolonged or extensive use can thin the skin or rarely cause effects throughout the body.

\medterm{Fluocinolone Acetonide} A corticosteroid medicine applied to the skin or scalp to reduce redness, swelling, and itching. Overuse can thin the skin, worsen infections, or rarely affect the body’s hormone system.

\medterm{Fluocinonide} A strong corticosteroid applied to selected inflammatory skin conditions to reduce itching, redness, and swelling. It should be used only as directed because prolonged or extensive use can damage skin or cause systemic effects.

\medterm{Fluocortolone} A corticosteroid used in selected skin or other inflammatory conditions. It reduces inflammation but can cause skin thinning, infection, or other steroid effects, especially when used for a long time or over large areas.

\medterm{Fluorescein} A bright fluorescent dye used in eye examinations and some imaging tests. In the eye, it highlights scratches, ulcers, and abnormal blood-vessel leakage; rare allergic reactions can occur after injection.

\medterm{Fluorescein Angiography} An eye test in which fluorescent dye is injected into a vein and photographs track its movement through retinal blood vessels. It can show leakage, blockage, or abnormal vessels in retinal disease.

\medterm{Fluorescein Eye Stain} A test in which dye is placed on the eye and viewed under blue light. It highlights scratches, ulcers, foreign-body injuries, and some dry-eye surface damage.

\medterm{Fluorescence} The emission of visible light by a substance after it absorbs light or other energy. Clinicians and researchers use it to identify cells, chemicals, microbes, or structures in laboratory and imaging tests.

\medterm{Fluorescence In Situ Hybridization} A laboratory test that uses glowing DNA probes to find specific chromosome regions or gene changes in cells. It can detect selected missing, extra, rearranged, or amplified genetic material.

\synonyms
Fluorescence In Situ Hybridization

\medterm{Fluorescence In-Situ Hybridization} A targeted genetic test using fluorescent DNA probes to locate selected chromosome sequences. It can identify some extra or missing chromosome material and gene rearrangements but does not examine the entire genome.

\medterm{Fluorescence Microscopy} A microscope method that uses fluorescent labels to make selected cells, molecules, or microbes visible. It helps examine tissue structure, cell proteins, genetic material, and infectious organisms.


\medterm{Fluorescent Dye} A substance that absorbs light of one color and emits light of another. It can label cells, tissues, molecules, or microbes so they can be detected during laboratory or imaging examinations.

\medterm{Fluoridation} Adding a controlled amount of fluoride to water or dental products to strengthen tooth enamel and reduce cavities. Excessive fluoride during tooth development can cause dental fluorosis, which appears as enamel discoloration or pits.

\medterm{Fluoride} A mineral that strengthens tooth enamel, helps repair early damage, and reduces acid production by cavity-causing bacteria. It is found in some water supplies, toothpaste, mouth rinses, and professional dental treatments.

\medterm{Fluoride In Diet} Fluoride obtained from drinking water and foods or drinks made with fluoridated water. The right amount helps protect teeth, while excessive intake, especially in children, can damage developing enamel.

\medterm{Fluoride Overdose} Excessive fluoride swallowed over a short time. It can cause nausea, vomiting, abdominal pain, low blood calcium, muscle spasms, abnormal heart rhythms, or seizures and requires urgent poison-control or medical advice.

\medterm{Fluoride Poisoning} Illness caused by swallowing too much fluoride. Early symptoms include nausea, vomiting, and abdominal pain; severe poisoning can lower blood calcium and cause muscle spasms, seizures, or dangerous heart rhythms.

\medterm{Fluoride Varnish} A concentrated fluoride coating painted onto teeth by a dental professional to prevent or slow cavities. It hardens quickly and is used at intervals based on age, tooth condition, and cavity risk.

\medterm{Fluorimetry} A laboratory method that measures light emitted by a fluorescent substance. It can detect or estimate very small amounts of chemicals, medicines, proteins, or biological material.

\medterm{Fluorine} A highly reactive chemical element found in compounds such as fluoride. Small controlled amounts of fluoride help protect teeth, while excessive exposure can poison the body or damage developing bones and teeth.

\medterm{Fluorine-18} A radioactive form of fluorine used to label tracers for positron-emission tomography scans. Its short half-life allows useful imaging while limiting how long radioactivity remains in the body.

\medterm{Fluorometholone} A corticosteroid eye medicine used for selected inflammatory conditions. Prolonged use can raise eye pressure, delay healing, worsen infection, or increase cataract risk, so it requires medical supervision.

\medterm{Fluorophotometry} A method that measures fluorescent light from the eye or another sample. It can help estimate substances or tissue properties, but results must be interpreted with the test method and clinical findings.

\medterm{Fluoropyrimidine} A class of anticancer medicines that resembles substances used to build DNA and RNA. They disrupt genetic-material production in rapidly dividing cells and can cause mouth sores, diarrhea, low blood counts, and hand-foot symptoms.

\medterm{Fluoroquinolone} An antibiotic that kills bacteria by blocking enzymes needed for bacterial DNA copying. It is reserved for selected infections because it can cause tendon injury, nerve symptoms, heart-rhythm changes, blood-sugar changes, or severe diarrhea.

\medterm{Fluoroquinolones} A group of antibiotics that kill bacteria by blocking enzymes required for bacterial DNA replication. They treat selected infections but are avoided when safer options work because of potentially serious tendon, nerve, heart, and blood-sugar effects.

\medterm{Fluoroscopy} An X-ray method that produces moving images inside the body. It helps guide procedures and examine swallowing, blood vessels, joints, and digestive or urinary passages; contrast dye may be used and radiation exposure occurs.

\medterm{Fluorosis} A change in developing tooth enamel caused by too much fluoride during childhood. It usually appears as white streaks or spots; more severe cases cause brown discoloration or pits.

\medterm{Fluorouracil} An anticancer medicine that interferes with DNA production. It is applied to selected precancerous skin lesions or given for some cancers; possible effects include mouth sores, diarrhea, low blood counts, and serious toxicity in susceptible people.

\medterm{Fluorouracil Side Effect} An unwanted effect of fluorouracil, such as mouth sores, diarrhea, nausea, low blood counts, hand-foot skin reactions, or heart problems. Severe symptoms require prompt medical review and treatment adjustment.

\medterm{Fluoxetine} An antidepressant that increases available serotonin in the brain. It treats depression and several other conditions, but may cause nausea, sleep or sexual problems, agitation, or rarely dangerous serotonin excess or suicidal thoughts.

\medterm{Fluoxymesterone} A synthetic androgen medicine used in selected hormone-related conditions. It can affect the liver, cholesterol, mood, fertility, prostate, and blood count and should be used only with medical supervision.

\medterm{Flupentixol} An antipsychotic medicine used in some countries for schizophrenia and other psychiatric conditions. It can reduce hallucinations and delusions but may cause movement problems, sleepiness, hormonal changes, or rare serious reactions.

\medterm{Fluphenazine} An antipsychotic medicine used to treat schizophrenia and related psychotic symptoms. It can cause stiffness, tremor, restlessness, sleepiness, hormonal changes, and rare but serious involuntary movements or fever.

\medterm{Fluphenazine Hydrochloride} The hydrochloride salt form of fluphenazine, an antipsychotic medicine for selected psychotic disorders. Monitoring is needed because it can cause movement disorders, sedation, hormonal effects, and other serious adverse reactions.

\medterm{Flurandrenolide} A corticosteroid applied to the skin to reduce inflammation, redness, and itching. Prolonged, extensive, or covered use can thin the skin, worsen infection, or allow steroid effects throughout the body.

\medterm{Flurazepam} A long-acting sedative medicine used for short-term treatment of insomnia in some settings. It can cause next-day drowsiness, falls, memory problems, dependence, and dangerous breathing suppression with alcohol or opioids.

\medterm{Flurbiprofen} A nonsteroidal anti-inflammatory medicine used for pain and swelling and, in some forms, eye inflammation. It can cause stomach bleeding, kidney problems, heart risks, or allergic reactions and should be used as directed.

\medterm{Flush} A sudden feeling of warmth or visible reddening caused by increased skin blood flow. In healthcare, flush can also mean washing a tube, body passage, or intravenous line with fluid.

\medterm{Flushable Reagent Stool Blood Test} Alternate index label; see \textit{Fecal occult blood test}.

\medterm{Flushing} Temporary redness and warmth of the skin caused by increased blood flow. Triggers include heat, exercise, emotion, alcohol, medicines, fever, or hormone disorders; flushing with wheezing, diarrhea, faintness, or palpitations needs evaluation.

\medterm{Flutamide} A medicine that blocks male-hormone effects and is used with other hormone treatment for prostate cancer. It can seriously injure the liver and may cause nausea, breast tenderness, or reduced blood counts.

\medterm{Flutter} A rapid, regular, repetitive movement or contraction. Atrial flutter is an abnormal heart rhythm caused by a fast electrical circuit in the upper chambers and can increase the risk of stroke.

\medterm{Flutter Valve} A handheld breathing device that creates gentle pressure and vibrations during exhalation. It can help loosen and move mucus in some people with bronchiectasis, cystic fibrosis, or chronic lung disease.

\medterm{Fluvastatin} A statin medicine that lowers low-density lipoprotein cholesterol by reducing cholesterol production in the liver. It lowers cardiovascular risk but can cause muscle symptoms, liver abnormalities, or rare muscle breakdown.

\medterm{Fluvoxamine} An antidepressant that increases serotonin signaling and is used mainly for obsessive-compulsive disorder and selected other conditions. It can cause nausea, sleep changes, sexual problems, medicine interactions, or rarely serotonin toxicity.

\medterm{Flux} An old, nonspecific term for diarrhea or an excessive discharge, especially from the bowel. Modern medical descriptions should identify the actual symptom, duration, severity, and suspected cause.

\medterm{FM} An abbreviation sometimes used for fibromyalgia, a long-term condition causing widespread pain, fatigue, poor sleep, and concentration problems. Abbreviations can be unclear, so the full term is preferred in patient records.

\synonyms
Fibromyalgia

\medterm{FMRI} Functional magnetic resonance imaging maps changes in brain blood flow and oxygen use while a person performs tasks or rests. It can help identify language or movement areas and does not use X-ray radiation.

\medterm{FMS} An abbreviation sometimes used for fibromyalgia syndrome, a long-term condition with widespread pain, fatigue, poor sleep, and cognitive symptoms. The full term is preferred because FMS can have other meanings.

\synonyms
Fibromyalgia

\medterm{FNA} An abbreviation for fine-needle aspiration, a procedure using a thin needle to collect cells or fluid from a lump or organ for laboratory examination. The full term is clearer in patient information.

\medterm{FND} Functional neurological disorder causes real weakness, abnormal movements, seizure-like episodes, sensory changes, or other nervous-system symptoms because nerve signaling is disrupted. Symptoms are not imagined; diagnosis is based on characteristic examination findings and may require tests to exclude other causes.

\synonyms
Conversion disorder; functional neurologic disorder

\medterm{Foam Cell} A cell containing many fat droplets, giving it a foamy appearance under a microscope. Foam cells, often immune cells that have taken up cholesterol, are common in early artery-wall plaque.

\medterm{Foam Cells} Cells filled with fat droplets, often immune cells that have taken up altered cholesterol. They commonly form early in artery-wall plaque and contribute to inflammation and atherosclerosis.

\medterm{Focal} Limited to one specific area rather than spread widely. A focal finding may involve an image, seizure, infection, injury, or neurologic problem and must be interpreted according to its location.

\medterm{Focal Adhesion} A specialized attachment site where a cell connects to its surrounding support material. It helps cells stick, sense physical forces, move, and send signals during growth and tissue repair.

\medterm{Focal Glomerulosclerosis} Scarring affecting some kidney filters or only parts of affected filters. It can cause protein to leak into urine, swelling, high blood pressure, and gradual loss of kidney function.
\medterm{Focal Neurologic Deficits} Loss of a particular function, such as strength, sensation, speech, vision, or coordination, because of a localized problem in the brain, spinal cord, or a nerve. Sudden onset may indicate stroke.

\medterm{Focal Neuropathy} Damage or irritation affecting one peripheral nerve or a small group of nerves. It may cause pain, tingling, numbness, weakness, or altered movement in the area supplied by those nerves.

\synonyms
Mononeuropathy

\medterm{Focal Nodular Hyperplasia} A usually harmless overgrowth of liver tissue related to an unusual local blood supply. It often causes no symptoms, is found on imaging, and rarely needs treatment unless its diagnosis is uncertain.

\medterm{Focal Sclerosis} Hardening or scarring limited to a particular area of tissue. In the kidney, it may describe scarring of some filters or parts of filters, which can impair filtration and cause protein in urine.

\medterm{Focal Segmental Glomerulosclerosis} A kidney disorder in which scars affect some filters or parts of them. It can cause heavy urine protein, swelling, high blood pressure, and kidney failure; causes include genetic changes, infections, medicines, and excess filtration.

\medterm{Focal Seizure} A seizure that begins in one area of the brain. It may cause unusual movement, sensation, behavior, emotion, or awareness and can spread to involve both sides of the brain.

\medterm{Focal-Onset Seizure} A seizure that begins in one area of the brain. Awareness may remain normal or become impaired, and symptoms depend on the affected area; the seizure may spread and cause whole-body stiffening and jerking.

\medterm{Focus Words} Words or short phrases used during relaxation or stress-management exercises to direct attention and support calm breathing. They may help coping but are not a treatment for an underlying medical disorder.

\medterm{FODMAP Diet} A short-term eating plan that reduces certain poorly absorbed carbohydrates that can cause gas, bloating, pain, or diarrhea. Foods are gradually reintroduced with professional guidance to identify personal triggers.

\medterm{FODMAPs} A group of short-chain carbohydrates that may be poorly absorbed in the small intestine and fermented by gut bacteria. In susceptible people they can cause gas, bloating, abdominal pain, or diarrhea.

\medterm{Foetal Acidosis} Excess acidity in fetal blood, usually caused by inadequate oxygen delivery or impaired gas exchange. Fetal heart-rate monitoring and, in selected situations, blood testing help assess its likelihood and severity.

\medterm{Foetal Alcohol Syndrome} A lifelong condition caused by alcohol exposure before birth, with characteristic facial features, poor growth, and brain-development problems affecting learning, behavior, or coordination. No amount of alcohol is known to be safe during pregnancy.

\medterm{Foetus} The developing human offspring from the end of the embryonic period, about nine weeks after fertilization, until birth. During this period, organs grow and mature and body functions become more coordinated.

\medterm{Folate} Vitamin B9, needed for DNA production, cell division, and healthy red blood cells. Adequate folate before and during early pregnancy lowers neural-tube-defect risk; high-dose supplements can hide vitamin B12 deficiency.

\medterm{Folate Antagonist} A medicine that blocks folate-related processes needed to make DNA. Depending on the drug, it may treat cancer, inflammatory disease, or bacterial infection, but can cause low blood counts, mouth sores, or liver problems.

\medterm{Folate Deficiency} Too little vitamin B9, caused by poor intake, increased need, malabsorption, alcohol use, or some medicines. It can cause large abnormal red blood cells, anemia, fatigue, sore tongue, and pregnancy complications.

\medterm{Folate Deficiency Anemia} Anemia caused by too little folate, or vitamin B9. Red blood cells become unusually large and immature, causing fatigue, weakness, pale skin, shortness of breath, or a sore tongue.

\synonyms
Vitamin B9 deficiency anemia

\medterm{Folate-Deficiency Anemia} Anemia caused by inadequate folate, poor absorption, increased requirements, alcohol use, or certain medicines. Blood tests help confirm the cause, and folate replacement treats the deficiency after vitamin B12 deficiency is considered.

\medterm{Foley Catheter} A flexible tube placed through the urethra into the bladder to drain urine. A small balloon holds it in place; risks include discomfort, blockage, urinary infection, and urethral injury.

\medterm{Folic Acid} A manufactured form of vitamin B9 needed for DNA production and healthy blood cells. Supplements prevent many fetal neural-tube defects when taken before conception and treat folate deficiency.

\medterm{Folic Acid - Test} A blood test measuring folate levels to investigate suspected deficiency or anemia with unusually large red blood cells. Results are interpreted with blood counts, vitamin B12 testing, diet, medicines, and symptoms.

\medterm{Folic Acid And Birth Defect Prevention} Taking folic acid before conception and during early pregnancy reduces the risk of neural-tube defects, which affect development of the fetal brain and spinal cord. Higher-risk people may need a prescribed dose.
\medterm{Folic Acid In Diet} Folate occurs naturally in leafy vegetables, beans, citrus fruits, and liver, while folic acid is added to some fortified foods. Intake supports blood formation and fetal development; needs increase during pregnancy.
\medterm{Folic Acid Supplementation} Taking vitamin B9 before conception and during early pregnancy to reduce neural-tube-defect risk. Dose depends on individual risk; supplementation also treats deficiency and supports DNA production and red-cell formation.

\medterm{Folinic Acid} An active form of folate used to treat or prevent folate-related toxicity and to support selected chemotherapy treatments. It can also treat certain folate deficiencies; its use depends on the clinical situation.

\medterm{Follicle} A small sac or organized tissue structure. A hair follicle produces a hair, while an ovarian follicle contains a developing egg; other follicles can occur in glands or lymphatic tissue.

\medterm{Follicle Stimulating Hormone} A hormone released by the pituitary gland. In females it helps ovarian follicles grow and produce estrogen; in males it supports sperm production. Levels vary with age, sex, and reproductive stage.

\medterm{Follicle-Stimulating Hormone} A pituitary hormone that supports egg-containing follicle development and estrogen production in females and sperm production in males.

\medterm{Follicle-Stimulating Hormone Blood Test} A blood test measuring a hormone involved in ovarian function and sperm production. It helps evaluate infertility, menstrual changes, puberty problems, menopause, testicular or ovarian failure, and pituitary disorders.

\medterm{Follicular} Relating to a follicle, especially a hair follicle or an ovarian follicle. A follicular skin problem is centered around a hair and may appear as a bump, pustule, plug, or inflamed area.

\medterm{Follicular Cyst} A usually harmless ovarian cyst that forms when a follicle does not release an egg or does not shrink normally. It often causes no symptoms and disappears within a few menstrual cycles without treatment.

\medterm{Follicular Fluid} The liquid inside an ovarian follicle surrounding a developing egg. It contains hormones and other substances that support follicle growth and help coordinate ovulation.

\medterm{Follicular Hyperkeratosis} Excess keratin blocking hair follicles and causing small rough bumps. It commonly occurs with keratosis pilaris and, less often, with nutritional deficiency or other skin conditions.

\medterm{Follicular Lymphoma} A usually slow-growing cancer of immune cells that often causes painless swelling of lymph nodes. It may involve bone marrow or other organs and can sometimes change into a faster-growing lymphoma.

\medterm{Follicular Thyroid Carcinoma} A thyroid cancer that begins in follicle-forming cells and can spread through the bloodstream, especially to bone or lungs. Tissue examination is needed to show invasion and distinguish it from a benign growth.

\medterm{Follicular Unit} A natural group of one or more hair follicles along with nearby skin glands, nerves, and tiny muscles. Hair-transplant procedures may move these units from a fuller area to a thinning area.

\medterm{Follicular Unit Transplantation} Hair-restoration surgery that moves natural groups of hair follicles from a donor area to thinning or bald areas. Results take months, and possible problems include scarring, infection, and continued hair loss.

\synonyms
FUT

\medterm{Folliculitis} Inflammation or infection around a hair follicle, causing small itchy, red, or tender bumps and sometimes pus. Causes include bacteria, yeast, friction, blocked skin, shaving, and ingrown hairs.

\synonyms
Barber's Itch


\medterm{Follow-Up Recommendations} A plan for repeat examination, testing, imaging, treatment, or specialist review after an initial finding. The plan depends on the diagnosis, symptoms, test results, personal risks, and accepted clinical guidance.

\medterm{Fomentation} Applying a warm, moist cloth or material to part of the body to soothe pain, relax muscles, or improve comfort. Excessive heat can burn skin and should be avoided.

\medterm{Fomepizole} An antidote for poisoning with methanol or ethylene glycol. It blocks formation of their toxic breakdown products and may be combined with intensive supportive care and dialysis.

\medterm{Fomite} An object or surface contaminated with an infectious organism that can transfer it to another person. Examples include towels, equipment, toys, or needles, although the importance of surface spread varies by organism.

\medterm{Fomites} Objects or surfaces that may carry infectious organisms and transfer them through contact. Examples include shared towels, medical equipment, toys, and other contaminated items; cleaning and hand hygiene reduce risk.

\medterm{Fondaparinux} An injectable anticoagulant that prevents formation or extension of blood clots by blocking activated factor X. It is used for selected clotting disorders and prevention after surgery; bleeding is its main risk.

\medterm{Fondaparinux Sodium} The sodium salt form of fondaparinux, an injectable medicine that reduces clotting by blocking activated factor X. It treats or prevents selected blood clots but can cause serious bleeding and requires dose adjustment in kidney disease.

\medterm{Fonofos} A highly poisonous organophosphate insecticide. Exposure can cause sweating, excess saliva, vomiting, small pupils, muscle twitching, breathing difficulty, seizures, or weakness and requires urgent poison-control or emergency care.

\medterm{Fonsecaea} A group of darkly pigmented fungi that can cause slow-growing skin and soft-tissue infections called chromoblastomycosis. Infection usually follows injury and may produce warty plaques or nodules, especially in tropical regions.

\medterm{Fontan Procedure} An operation for some children born with only one effective pumping chamber. It directs blood returning from the body to the lungs without a separate pumping chamber, improving circulation but requiring lifelong follow-up for rhythm, heart, liver, clotting, and protein-loss problems.

\synonyms
Fontan Operation; Fontan Circulation; Total Cavopulmonary Connection (TCPC)

\medterm{Fontanel} A soft gap between an infant’s skull bones that allows the brain and skull to grow and helps the skull mold during birth. Its size and firmness can provide clues about hydration or pressure inside the skull.

\medterm{Fontanel (Fontanelle)} A soft membranous space between an infant’s skull bones. It normally changes as the skull grows; a persistently bulging or markedly sunken soft spot can signal illness and needs assessment.

\medterm{Fontanelle} A soft membranous space between an infant’s skull bones. Fontanelles allow growth and molding during birth; the front one usually closes during the second year, while the back one closes earlier.

\medterm{Fontanelles - Bulging} A tense soft spot that remains outwardly bulging when an infant is calm and upright. It may indicate increased pressure inside the skull, infection, bleeding, or another serious problem requiring urgent assessment.

\medterm{Fontanelles - Enlarged} Soft spots on an infant’s skull that are larger than expected for age. Causes include normal variation, prematurity, thyroid disease, vitamin D deficiency, bone disorders, or increased pressure inside the skull.

\medterm{Fontanelles - Sunken} A noticeably depressed soft spot in an infant who is calm and upright. It commonly suggests dehydration, especially with poor feeding, fewer wet diapers, vomiting, or diarrhea, and needs medical assessment.

\medterm{Food Additives} Substances added to food to preserve it or change its flavor, color, texture, or stability. Most are safe in approved amounts, but some people may have intolerance or allergy to particular additives.

\medterm{Food Allergies} Immune reactions to particular food proteins. They may cause hives, swelling, vomiting, wheezing, or life-threatening anaphylaxis; strict avoidance and an emergency plan are important for people with severe reactions.
\synonyms
Food allergy

\medterm{Food Allergy} An immune reaction to a food protein. It may cause hives, swelling, vomiting, abdominal symptoms, wheezing, or anaphylaxis. Common triggers include milk, egg, peanut, tree nuts, fish, shellfish, soy, and wheat. Severe reactions require an emergency plan.

\synonyms
Allergy Food (Food Allergy)
Allergy, Food

\medterm{Food Allergy, Egg} An immune reaction to egg proteins, most common in children. It can cause hives, vomiting, breathing problems, or anaphylaxis; evaluation guides avoidance, emergency planning, and whether the allergy has resolved.

\synonyms
Egg allergy

\medterm{Food Allergy, Milk} An immune reaction to milk proteins that can cause eczema, hives, vomiting, diarrhea, wheezing, or anaphylaxis. It is common in young children, and many outgrow it; safe nutritional substitutes are important.

\synonyms
Milk allergy

\medterm{Food Consumption} The amount, type, and pattern of food and drink a person takes in. It is influenced by appetite, culture, access, habits, health, activity, medicines, and the body’s nutritional needs.

\medterm{Food Hygiene} Practices that prevent food contamination and foodborne illness, including washing hands, separating raw and cooked foods, cooking thoroughly, chilling promptly, and using safe water.

\medterm{Food Intolerance} A troublesome reaction to a food that does not involve the immune system. Lactose intolerance is common; symptoms may include gas, bloating, cramps, or diarrhea and depend on the food and amount.

\medterm{Food Jags} A child’s repeated preference for only a few foods or refusal of foods previously accepted. It is common during development but needs assessment if it limits nutrition, growth, social eating, or causes distress.

\medterm{Food Labeling} Information on packaged food showing ingredients, allergens, serving size, nutrients, calories, dates, and health claims. Reading labels helps people manage allergies, medical diets, and overall nutrition.

\medterm{Food Poisoning} Illness from eating food contaminated with germs, toxins, or harmful chemicals. Nausea, vomiting, diarrhea, abdominal cramps, and fever are common.

\medterm{Food Poisoning (Foodborne Illness)} Illness caused by eating or drinking material contaminated with germs, toxins, or harmful chemicals. Symptoms commonly affect the digestive system, but some infections can also damage the nervous system or other organs.

\medterm{Food Poisoning Prevention} Reduce risk by washing hands, separating raw foods, cooking thoroughly, refrigerating promptly, avoiding unsafe water, and following food recalls. High-risk people should avoid foods commonly linked to serious infection.

\medterm{Food Poisoning, Campylobacter} A foodborne infection caused by Campylobacter bacteria, often from undercooked poultry, unpasteurized milk, or contaminated water. It causes diarrhea, cramps, and fever and may rarely trigger bloodstream infection or nerve inflammation.

\medterm{Food Safety} Handling, storing, preparing, and labeling food in ways that prevent germs, toxins, harmful chemicals, and allergens from causing illness.

\medterm{Food Supplement} A product taken to add nutrients or other substances to the diet, such as vitamins, minerals, herbs, or protein. Supplements can interact with medicines and are not automatically safe or effective.

\medterm{Food-Borne Disease} Illness acquired from food or water contaminated with infectious organisms, toxins, or harmful chemicals. It may cause digestive symptoms or affect other organs depending on the contaminant.

\medterm{Food-Borne Illness} Illness caused by eating food contaminated with bacteria, viruses, parasites, toxins, or chemicals. Symptoms commonly include nausea, vomiting, diarrhea, cramps, and fever; safe handling and cooking help prevent it.

\synonyms
Food poisoning; foodborne illness

\medterm{Foodborne Disease} Illness acquired by eating or drinking material contaminated with germs, toxins, or harmful chemicals. Severity ranges from brief diarrhea to dehydration, organ damage, or life-threatening infection.

\medterm{Foodborne Illness} Illness resulting from contaminated food or water, including infections and poisoning from germs, toxins, or chemicals. Serious disease is more likely in pregnancy, infancy, older age, or weakened immunity.


\medterm{Foot} The part of the lower limb below the ankle, made of bones, joints, muscles, tendons, ligaments, nerves, and blood vessels. Its shape and arches support standing, balance, shock absorption, and walking.



\medterm{Foot Drop} Difficulty lifting the front of the foot because of weakness in the ankle and foot muscles. It may cause tripping or a high-stepping walk and can result from nerve injury, spinal disease, muscle disease, or stroke.

\synonyms
Drop Foot

\medterm{Foot Fracture} A break in one or more foot bones caused by sudden injury, repeated stress, or weakened bone. Pain, swelling, bruising, and difficulty bearing weight are common; treatment ranges from protection and immobilization to surgery.

\synonyms
Broken foot

\medterm{Footling Breech} A fetal position in which one or both feet point downward toward the birth canal. It increases the risk of umbilical-cord prolapse and difficulty delivering the baby, so delivery planning requires specialist obstetric assessment.

\synonyms
Incomplete Breech; Footling Breech Presentation

\medterm{Foot Pain} Pain in the foot caused by injury, overuse, arthritis, nerve problems, infection, or poor circulation. Persistent pain, inability to bear weight, numbness, discoloration, fever, or an open wound requires assessment.


\medterm{Foot, Leg, And Ankle Swelling} Enlargement of the lower limb from fluid, inflammation, injury, infection, vein disease, medicines, or heart, kidney, or liver problems. Sudden one-sided swelling, pain, redness, or breathlessness needs urgent assessment.

\medterm{Foramen} A natural opening or passage through a bone or other body structure, usually allowing a nerve, blood vessel, or other tissue to pass.

\medterm{Foramen Ovale} An opening between the upper chambers of the fetal heart that directs blood away from the unused lungs. It usually closes after birth; persistence is called a patent foramen ovale.

\medterm{Foraminotomy} Surgery that enlarges an opening where a spinal nerve leaves the backbone, relieving pressure caused by arthritis, a disc problem, or other narrowing. Risks include infection, bleeding, nerve injury, and persistent symptoms.

\medterm{Force Plates} Platforms that measure forces beneath the feet during standing, walking, or other movement. Clinicians and researchers use them to assess balance, gait, symmetry, and how weight is distributed.

\medterm{Forced Expiratory Volume} The amount of air forcefully exhaled during a measured time after taking a full breath. The first-second value helps assess narrowed or restricted airways and is interpreted with other lung-test results.

\medterm{Forced Oscillation Technique} A breathing test in which a device sends gentle pressure waves through the airways while a person breathes normally. It measures how easily air moves and can help assess lung disease when forceful breathing tests are difficult.

\medterm{Forced Vital Capacity} The total amount of air forcefully exhaled after a person breathes in fully. A low value may occur when the lungs are small or stiff or when airways prevent complete emptying.

\medterm{Forceps} A handheld instrument used to grasp, hold, or move tissue or medical materials. Different designs have smooth or toothed tips and are chosen according to the procedure.

\medterm{Forceps (Obstetric)} Paired curved instruments placed around the fetal head during selected assisted vaginal births. They can help complete delivery when the baby or pregnant person needs delivery sooner and conditions are suitable.

\synonyms
Obstetric forceps; delivery forceps

\medterm{Forceps (Surgical)} Handheld instruments used during surgery to hold, grasp, move, or manipulate tissue, sutures, dressings, or other materials. Their shape and pressure are selected to limit tissue injury.

\synonyms
Surgical grasping instrument

\medterm{Forceps (Vascular)} Surgical forceps designed to handle blood vessels, clots, tissue, or medical materials during vascular procedures. Appropriate use helps avoid crushing or tearing delicate vessels.

\synonyms
Vascular surgical forceps

\medterm{Forceps Delivery} An assisted vaginal birth in which forceps guide the baby’s head out of the birth canal. It may be used for prolonged labor, concerning fetal findings, or inability to push safely; tears and newborn injury are possible.

\medterm{Forceps-Assisted Delivery} A vaginal delivery assisted by placing forceps around the fetal head. It is considered only when the cervix is fully open, the head’s position is known, and vaginal birth is judged safe; maternal and newborn injuries can occur.

\medterm{Forcible Intercourse} Sexual intercourse without freely given consent, involving force, threats, coercion, or inability to consent. It is sexual violence; survivors may need immediate safety support, medical care, emergency contraception, infection prevention, and counseling.

\medterm{Fordyce'S Angiokeratomas} Small, usually harmless dark-red or purple spots caused by widened tiny blood vessels, commonly on the scrotum or vulva. They may bleed after injury; unusual or changing lesions should be examined.

\medterm{Forearm} The part of the arm between the elbow and wrist, containing the radius and ulna bones, muscles, tendons, nerves, and blood vessels.

\medterm{Forearm Circumference} The distance around the forearm at a specified point. Measurement can help monitor growth, muscle loss, swelling, nutrition, or changes during treatment and rehabilitation.

\medterm{Forehead} The part of the face above the eyebrows and in front of the frontal skull bone. It contains skin, muscles, nerves, and blood vessels and helps with facial expression and protection.

\medterm{Forehead Lift} Cosmetic surgery that raises drooping eyebrows or forehead skin and may reduce wrinkles. Risks include scarring, numbness, asymmetry, hairline changes, bleeding, infection, and nerve injury.

\synonyms
Brow lift

\medterm{Foreign Body} An object that is not normally part of the body and becomes lodged in tissue or a body passage, such as glass, metal, food, or a splinter. It may cause injury, blockage, infection, or inflammation.

\medterm{Foreign Body Airway Obstruction} Partial or complete blockage of the breathing passages by an inhaled object or material. Severe choking with inability to speak, cough, or breathe is life-threatening and requires immediate first aid and emergency help.

\medterm{Foreign Body Aspiration} Inhalation of an object or food into the airway, especially in young children. Choking, persistent cough, wheezing, noisy breathing, or reduced breath sounds may occur; complete blockage is an emergency.

\medterm{Foreign Body In The Nose} An object lodged in a nasal passage, often in a child. One-sided blockage, foul discharge, bleeding, pain, or difficulty breathing may occur; safe removal prevents tissue damage and infection.

\medterm{Foreign Body Reaction} Inflammation around material that remains in tissue, such as a splinter, stitch, implant, or device. It may cause a lump, redness, pain, drainage, or scar tissue and can resemble infection.

\medterm{Foreign Object - Inhaled} An object or material that enters the airway during breathing. Coughing, choking, wheezing, noisy breathing, or one-sided reduced breath sounds may occur; breathing difficulty requires emergency care.

\medterm{Foreign Object - Swallowed} An object that is accidentally swallowed and enters the digestive tract. Many small smooth objects pass naturally, but batteries, sharp objects, magnets, or objects causing pain or blockage require urgent medical assessment.

\medterm{Foreign-Body Aspiration} Entry of an object or food into the voice box, windpipe, or bronchial tubes, especially in young children. Sudden choking, cough, wheezing, or noisy breathing may occur; complete blockage is an immediate emergency.

\medterm{Forensic} Relating to the use of scientific or medical knowledge in legal investigations, court proceedings, identification, injury assessment, or determining the cause and manner of death.

\medterm{Forensic Anthropology} The examination of human remains, especially bones, for legal investigations. It can help determine whether remains are human, estimate biological characteristics, assess injury, and support identification.


\medterm{Forensic Evidence} Physical, biological, digital, or documentary material collected and analyzed during a legal investigation. Examples include blood, DNA, fibers, fingerprints, toxicology samples, photographs, and medical records.

\medterm{Forensic Examination} A structured medical or scientific assessment performed for legal purposes. It may document injuries, collect evidence after assault, assess impairment, examine a death, or analyze a scene or specimen.

\medterm{Forensic Genetics} The use of DNA testing to identify people, establish biological relationships, or compare biological evidence in legal investigations. Results require proper collection, chain of custody, and expert interpretation.

\medterm{Forensic Odontology} The use of dental knowledge in legal investigations. Comparing teeth, dental records, and radiographs can help identify remains and document dental injuries or other evidence.

\medterm{Forensic Pathology} The medical study of deaths and injuries for legal purposes. Forensic pathologists investigate sudden, unexpected, violent, or suspicious deaths and help determine the cause and manner of death.

\medterm{Forensic Psychiatry} The application of psychiatry to legal questions, such as ability to stand trial, mental state during an alleged offense, risk assessment, involuntary treatment, and decision-making capacity.

\medterm{Forensic Radiography} The use of X-rays, computed tomography, magnetic resonance imaging, or other imaging in legal investigations. It can identify people, document injuries, locate objects, assess trauma, and support death investigations.

\medterm{Forensic Toxicology} The testing of blood, urine, hair, tissue, or other specimens for medicines, alcohol, drugs, poisons, and their breakdown products to assess possible roles in illness, impairment, poisoning, or death.

\medterm{Forensic Toxicology Screen} An initial laboratory test of biological samples for medicines, alcohol, drugs, poisons, or their breakdown products. A positive screen usually needs confirmation and measurement before legal or medical conclusions are made.

\medterm{Foreskin} The movable fold of skin covering the head of the penis. It protects the sensitive surface and contains nerves; it may be removed by circumcision.

\medterm{Foreskin (Prepuce)} The fold of skin covering and protecting the head of the penis. It normally becomes retractable over time; forceful retraction can cause pain, tearing, or scarring.

\synonyms
Prepuce

\medterm{Formaldehyde} A strong-smelling chemical used in manufacturing and tissue preservation. It can irritate the eyes, skin, and airways, cause allergic reactions, and increase cancer risk with substantial or long-term exposure.

\medterm{Formalin} A water-based solution of formaldehyde used to preserve tissue for laboratory examination. It can irritate or burn the skin, eyes, and airways and must be handled with appropriate ventilation and protection.

\medterm{Formalin-Fixed Paraffin-Embedded Tissue} Tissue preserved in formalin, processed, and embedded in wax so thin sections can be examined. These blocks support microscopic, protein, and many genetic tests, although preservation can alter some results.

\medterm{Formate} A salt or ester of formic acid. In medicine, formate is important because it is a toxic breakdown product of methanol and can contribute to severe acidosis and vision loss after poisoning.

\medterm{Formed Elements Of Blood} The cells and cell fragments suspended in blood: red blood cells carry oxygen, white blood cells help defend against infection, and platelets help form clots and stop bleeding.

\medterm{Formerly Churg-Strauss Syndrome} An older name for eosinophilic granulomatosis with polyangiitis, an inflammatory blood-vessel disease associated with asthma and high eosinophil levels. It can damage the lungs, nerves, skin, heart, kidneys, or digestive organs.
\medterm{Formetanate} A carbamate insecticide that can poison the nervous system by disrupting acetylcholine breakdown. Exposure may cause sweating, vomiting, small pupils, muscle weakness, breathing difficulty, or seizures and requires urgent medical advice.

\medterm{Formication} The feeling that insects are crawling on or under the skin when none are present. It may occur with nerve disease, menopause, substance withdrawal, medicines, or psychiatric illness and should be medically evaluated if persistent.

\medterm{Formoterol Fumarate} A long-acting inhaled medicine that relaxes airway muscles and helps prevent asthma or chronic obstructive lung disease symptoms. It is not a rescue medicine for sudden severe breathing difficulty and is often combined with an inhaled corticosteroid.

\medterm{Formothion} An organophosphate insecticide that can poison the nervous system. Exposure may cause sweating, excess saliva, vomiting, muscle twitching, breathing difficulty, seizures, or weakness and requires urgent emergency or poison-control advice.

\medterm{Formparanate} A pesticide-related chemical whose health effects depend on the formulation, dose, and route of exposure. Suspected poisoning should be treated as urgent and assessed through poison-control or emergency services.

\medterm{Formula Feed} Prepared infant nutrition given by bottle, cup, or feeding tube when breast milk is unavailable, insufficient, or not chosen. The formula and preparation method should match the infant’s age and medical needs.

\medterm{Formula Feeding} Feeding an infant with commercially prepared formula instead of or in addition to breast milk. Safe preparation, clean equipment, correct dilution, and age-appropriate amounts help prevent infection, dehydration, and poor nutrition.

\medterm{Formula, Infant} Commercially prepared nutrition for infants, available in standard, soy-based, extensively hydrolyzed, amino-acid, or other specialized forms. Choice depends on age, tolerance, growth, and medical needs; preparation must follow instructions exactly.

\medterm{Formulae} The plural of formula. In healthcare, the word may mean mathematical equations, medicine formulations, or prepared nutritional products such as infant formula.

\medterm{Formulary} A list of medicines approved or preferred by a health system, insurer, or institution. It may show coverage, restrictions, alternatives, and requirements for approval, but does not replace a clinician’s prescribing judgment.

\medterm{Fornix} A recess or arch-shaped structure. In the eye, the conjunctival fornix is the fold where the eyelid lining joins the surface covering the eyeball; the term also describes parts of the brain and vagina.

\medterm{Forrester Classification} A system that groups people with severe heart failure or shock according to blood flow and congestion. “Warm” means adequate flow and “cold” reduced flow; “dry” means little congestion and “wet” significant congestion, helping guide urgent treatment.

\synonyms
Forrester Hemodynamic Subsets; Forrester Score

\medterm{Fortification} Adding vitamins or minerals to foods to improve nutritional intake and prevent deficiency. Examples include iodine in salt, vitamin D in some dairy products, and folic acid in many grain products.

\medterm{Fortified Wine} Wine with added distilled alcohol, giving it a higher alcohol content than ordinary wine. Excess use can cause intoxication, dependence, liver disease, injuries, and interactions with medicines.


\medterm{Fosaprepitant Dimeglumine} An intravenous medicine converted to aprepitant that blocks a brain signaling pathway involved in nausea and vomiting. It helps prevent chemotherapy-related sickness and can interact with other medicines.

\medterm{Foscarnet Sodium} An antiviral medicine used for selected serious or resistant cytomegalovirus and herpesvirus infections. It can damage the kidneys and disturb blood salts, so hydration and laboratory monitoring are important.

\medterm{Fosfomycin} An antibiotic that blocks an early step in bacterial cell-wall production. It is used for selected urinary or difficult-to-treat infections; diarrhea, allergic reactions, and resistance can occur.

\medterm{Fosinopril} An angiotensin-converting enzyme inhibitor used for high blood pressure and some forms of heart failure. It relaxes blood vessels and reduces fluid strain; cough, high potassium, kidney problems, or swelling can occur.

\medterm{Fossa} A shallow hollow or depression in a bone or other body structure. It may provide space for a joint, muscle, blood vessel, nerve, gland, or part of an organ.

\medterm{Fossula} A small pit or shallow depression in an anatomical structure. Examples include small depressions in a bone or on a tooth; the term describes shape and location rather than a disease.


\medterm{Fotemustine} A chemotherapy medicine that can cross into the brain and is used in some countries for selected advanced cancers. It can suppress bone marrow, increasing the risk of infection and bleeding.

\medterm{Found Dead} A circumstance in which a person is discovered deceased without the death having been observed. Medical or forensic examination may determine identity, timing, cause, and manner of death.

\medterm{Four-Dimensional Ultrasound} Three-dimensional ultrasound shown as movement over time, creating a live view. It may help assess fetal anatomy or other structures, but image quality and diagnostic value depend on the examination and equipment.

\medterm{Fourchette} The small fold of tissue at the back of the opening of the vulva, where the inner labial folds meet. It may stretch or tear during childbirth or sexual injury.

\medterm{Fourier Analysis} A mathematical method that separates a complex signal into simpler repeating frequencies. It is used in medical imaging, heart and brain signal analysis, and other measurements to identify patterns.

\medterm{Fournier Gangrene} A rapidly spreading, life-threatening infection that destroys tissue around the genitals, perineum, or anus. Severe pain, swelling, fever, skin discoloration, or feeling very ill requires immediate surgery and antibiotics.


\medterm{Fourth Finger} The ring finger, located between the middle finger and little finger. It contains bones, joints, tendons, nerves, blood vessels, and skin that support movement and grasping.

\medterm{Fourth Stage Of Labor} The first hours after the placenta is delivered, when the birthing person is closely observed for heavy bleeding, abnormal blood pressure, poor uterine contraction, pain, and other early complications.

\medterm{Fourth Venereal Disease} An old name for lymphogranuloma venereum, a sexually transmitted infection caused by particular Chlamydia strains. It may cause a small genital sore followed by painful groin swelling and, untreated, rectal complications.

\medterm{Fourth-Degree Tear} A severe childbirth tear extending from the vaginal area through the anal sphincter and into the rectal lining. It needs careful repair, antibiotics when appropriate, follow-up, and monitoring for bowel-control problems.

\medterm{Fovea} A tiny central depression in the retina responsible for the sharpest detailed and color vision. Damage to it can impair reading, face recognition, and central vision.

\medterm{Fovea Centralis} The central pit of the retina with the greatest concentration of cone cells, providing the sharpest central and color vision. It has no major blood vessels crossing directly over its center.

\medterm{Fowler Position} A semi-upright position with the head and upper body raised, commonly about 45 to 60 degrees. It may improve comfort and breathing or be used during examinations and procedures.

\medterm{Foxglove Poisoning} Poisoning from foxglove plants or medicines containing cardiac glycosides. It can cause nausea, vomiting, confusion, visual changes, dangerously slow or irregular heart rhythms, or cardiac arrest and requires emergency care.

\medterm{Fractional Excretion Of Sodium} The percentage of filtered sodium leaving the body in urine. It can help assess kidney function and distinguish reduced kidney blood flow from some tubular injuries, but diuretics and other illnesses can make it misleading.

\medterm{Fractional Excretion Of Urea (FEUrea)} The percentage of filtered urea excreted in urine. It can support assessment of kidney blood flow when diuretics affect sodium testing, but infection, bleeding, nutrition, and chronic kidney disease also influence the result.

\synonyms
FEUrea

\medterm{Fractional Flow Reserve} A pressure measurement made during coronary angiography to determine whether a narrowed heart artery significantly limits blood flow. The result helps guide decisions about stents or medical treatment.

\medterm{Fractional IPV} A smaller-than-usual dose of inactivated polio vaccine given into the skin rather than the muscle. When used in an approved schedule, it helps protect against polio; the number of doses and timing follow local immunization guidance.

\medterm{Fracture} A break or crack in a bone caused by injury, repeated stress, or weakened bone. Pain, swelling, bruising, deformity, or inability to use the part may occur; treatment restores alignment and supports healing.

\medterm{Fracture Displacement} Movement of broken bone pieces out of their normal alignment, including sideways shift, angling, shortening, or rotation. The need for reduction or surgery depends on the bone, joint involvement, age, and function.

\medterm{Fracture Fixation} Stabilization of a broken bone with a cast, splint, plate, screw, nail, wire, or external frame. The goal is to maintain alignment while the bone heals and, when possible, preserve movement.

\medterm{Fracture Liaison Service} A coordinated healthcare program for people with fragility fractures. It evaluates bone strength and underlying causes, arranges treatment and fall prevention, and follows patients to reduce future fractures.

\medterm{Fracture Of The Temporal Bone} A break in a skull bone containing parts of the ear and nearby nerves and blood vessels. It may cause hearing loss, dizziness, facial weakness, bleeding, or leakage of brain fluid and needs urgent assessment.

\medterm{Fracture Prevention} Measures that lower the chance of broken bones, including treating osteoporosis, improving strength and balance, preventing falls, correcting vitamin or mineral deficiency, and reviewing medicines that increase risk.

\medterm{Fracture, Arm} A break in the humerus, radius, or ulna caused by injury, a fall, or weakened bone. Pain, swelling, deformity, and loss of movement may occur; treatment ranges from immobilization to surgery.

\synonyms
Broken arm

\medterm{Fracture, Growth Plate} A break involving the cartilage growth area of a child’s or adolescent’s bone. It can disturb future growth if displaced or untreated, so prompt imaging and specialist management may be needed.

\synonyms
Growth plate fractures; physeal fractures

\medterm{Fracture, Hip} A break in the upper femur near the hip joint, often after a fall in older adults or major trauma in younger people. It causes pain and inability to walk and commonly needs urgent surgery.

\synonyms
Hip fracture

\medterm{Fracture, Leg} A break in the tibia, fibula, or both lower-leg bones, usually from trauma, a fall, or sports injury. Pain, swelling, deformity, and inability to bear weight may occur; treatment depends on stability and alignment.

\synonyms
Broken leg

\medterm{Fracture, Stress} A small crack caused by repeated loading or overuse rather than one major injury. It commonly affects the foot or lower leg and causes pain that worsens with activity; continued stress can complete the break.

\synonyms
Stress fractures; fatigue fractures

\medterm{Fracture-Fixation Device} A plate, screw, rod, pin, wire, or external frame used to hold broken bone pieces in alignment during healing. Choice depends on the fracture, bone quality, soft-tissue injury, and expected movement.

\medterm{Fractured Clavicle In The Newborn} A broken collarbone in a newborn, usually related to birth stress. The baby may move one arm less or cry when handled; most heal well with gentle support and follow-up.

\medterm{Fractured Elbow} A break involving one of the bones forming the elbow. It can cause pain, swelling, deformity, and limited movement and may injure nearby nerves or blood vessels.

\medterm{Fractured Ilium} A break in the broad upper part of the pelvic bone, usually from major trauma or weakened bone. It may cause severe pain and can accompany injury to pelvic organs or blood vessels.

\medterm{Fractured Nose} A break in the nasal bones or supporting cartilage, usually after a blow. Pain, swelling, bleeding, and deformity may occur; a collection of blood inside the septum requires urgent treatment.

\synonyms
Broken nose; nasal fracture

\medterm{Fractured Rib} A break in one or more ribs, usually from injury or forceful coughing. Sharp pain worsens with breathing or movement; lung collapse, bleeding, or injury to nearby organs must be considered after significant trauma.

\synonyms
Broken ribs; rib fracture

\medterm{Fragile X Syndrome} An inherited condition caused by a change in the FMR1 gene. It can affect learning, speech, behavior, social development, and physical features and is a common inherited cause of intellectual disability.

\medterm{FRAIL Scale} A brief questionnaire assessing fatigue, ability to climb stairs, walking ability, illness burden, and unintentional weight loss. It screens for frailty risk but does not replace a complete medical and functional assessment.

\medterm{Frailty} A state of reduced physical reserve that makes an older or seriously ill person more vulnerable to falls, infection, hospitalization, disability, and other stress. Signs include weakness, slow walking, exhaustion, inactivity, and unintentional weight loss.

\medterm{Framboesia} Another name for yaws, a chronic infection caused by Treponema pallidum pertenue. It spreads mainly through skin contact in tropical areas and can cause skin lesions, bone pain, and deformity if untreated.

\medterm{Frameshift Mutation} A genetic change caused by inserting or deleting DNA bases in numbers that are not multiples of three. This shifts how the genetic code is read and often produces a severely altered or nonworking protein.

\medterm{Framycetin} An aminoglycoside antibiotic applied to selected bacterial skin, ear, or eye infections. It can cause local allergy or irritation; prolonged or extensive use may allow absorption and antibiotic-related toxicity.

\medterm{Francisella} A group of small bacteria that includes Francisella tularensis, which causes tularemia. People can acquire infection through ticks, animals, contaminated water, food, or inhaled particles.

\medterm{Francisella Tularensis} A bacterium that causes tularemia, usually acquired through ticks, deer flies, infected animals, contaminated water or food, or inhalation. It can cause fever, an ulcer with swollen lymph nodes, pneumonia, or severe systemic illness.

\medterm{Francisellaceae} A family of small bacteria that includes Francisella species. Some cause tularemia or related infections in animals and people; disease type depends on the species and route of exposure.

\medterm{Frank Breech} A fetal position in which the buttocks are nearest the birth canal, the hips are bent, and the knees are straight with the feet near the head. Delivery planning depends on gestational age, fetal size, and clinical circumstances.

\synonyms
Extended Breech; Frank Breech Presentation

\medterm{Frank-Starling Mechanism} The heart normally contracts more strongly when its chambers fill with more blood, up to a limit. This helps match the amount pumped with the amount returning to the heart.

\medterm{Frasier Syndrome} A rare inherited disorder caused by changes in the WT1 gene. It can cause abnormal sexual development, progressive kidney disease, and increased risk of tumors in the gonads.

\medterm{Fraternal Twins} Twins formed when two separate eggs are fertilized by two separate sperm. They are genetically like ordinary siblings and may be the same or different sex.

\medterm{FRAX Tool} A calculator estimating a person’s ten-year probability of a major osteoporosis-related fracture or hip fracture. It uses factors such as age, sex, bone density, previous fracture, and other risks and supports treatment decisions.
\synonyms
FRAX fracture-risk calculator

\medterm{Freckle} A small, flat, tan or light-brown area of increased skin pigment that often darkens with sunlight. Freckles are usually harmless, but a changing, bleeding, irregular, or unusual spot needs examination.

\medterm{Freckles} Small, flat pigmented spots that often become darker after sun exposure and reflect inherited skin characteristics. They are usually harmless; changes in size, shape, color, symptoms, or bleeding require assessment.

\medterm{Freckles (Ephelides)} Small, flat, light-brown skin spots caused by increased pigment production. They commonly darken with sunlight and fade when exposure decreases and are usually harmless.

\synonyms
Ephelides

\medterm{Free Fragment} A piece of spinal-disc material that has broken away from the main disc and may press on a nearby nerve. It can cause back or leg pain, numbness, weakness, or bladder and bowel problems.

\synonyms
Sequestered disk fragment

\medterm{Free Radical} A highly reactive molecule with an unpaired electron. The body makes some during normal metabolism, while excessive amounts can damage cell membranes, proteins, and DNA as part of oxidative stress.

\medterm{Free Radical Injury} Cell damage caused by excessive reactive molecules that alter fats, proteins, or DNA. It can contribute to inflammation and tissue injury, although the importance varies with the disease and exposure.

\medterm{Free T4 Test} A blood test measuring unbound thyroxine, the thyroid hormone available to tissues. It is usually interpreted with thyroid-stimulating hormone to assess overactive or underactive thyroid function and treatment response.

\medterm{Free Thyroxine Index} A calculated estimate of the thyroid hormone available to tissues, using total thyroxine and a measure of hormone binding. It can help assess thyroid function when binding-protein levels affect total hormone results.

\medterm{Free Water Clearance} A kidney measurement showing whether urine contains more or less water than blood plasma. Positive values indicate excretion of dilute urine; negative values indicate water retention and concentrated urine.

\medterm{Freeze Drying} A preservation process that freezes material and removes ice under low pressure. It helps preserve medicines, vaccines, biological samples, and foods while limiting damage from heat and moisture.


\medterm{Freezing} Exposure to very low temperature that can freeze tissue, cause frostbite, or dangerously lower body temperature. The term also describes preserving food, medicines, or samples by storing them at low temperature.

\medterm{Fremitus} A vibration felt or heard through body tissues. During a chest examination, transmitted voice vibrations can change with fluid, air, or solid tissue in or around the lungs.

\medterm{French Disease} An old name for syphilis, a sexually transmitted infection caused by Treponema pallidum. Untreated infection can progress from sores to widespread disease and later damage the brain, nerves, heart, or other organs.

\medterm{Frenectomy} A procedure that removes or releases a frenulum, a fold of tissue that may restrict tongue movement, pull on the gums, separate teeth, or interfere with dentures. Bleeding, infection, and scarring are possible.

\medterm{Frenulum} A small fold or band of skin or mucous membrane that anchors and limits movement of a body part, such as beneath the tongue, inside the lips, or on the penis.

\medterm{Frenulum Breve} An abnormally short penile frenulum that may cause pain, tearing, or bleeding during erection or sexual activity. Treatment depends on symptoms and may include observation or a minor surgical procedure.

\medterm{Frenum} A narrow fold of tissue that attaches one body surface to another and limits movement. Examples include the folds beneath the tongue, inside the lips, and beneath the foreskin.

\medterm{Frenum (Frenulum)} A fold of skin or mucous membrane that anchors a body part and limits its movement. Examples include the tongue fold, lip folds, and penile frenulum.

\synonyms
Frenulum

\medterm{Frequency} Urinating more often than usual, sometimes with only small amounts passed. It differs from producing a large total urine volume and may result from infection, pregnancy, diabetes, medicines, or bladder disorders.

\medterm{Frequency Of Urination} Needing to urinate more often than usual. Causes include increased fluid intake, urinary infection, pregnancy, diabetes, diuretic medicines, bladder irritation, prostate enlargement, or other urinary disorders.

\medterm{Frequency Range} The lowest to highest sound or signal frequencies that a device can detect, produce, or amplify effectively. In hearing care, it helps describe which pitches a hearing aid can process.

\medterm{Frequency Response} How strongly a device or body system responds to different sound or signal frequencies. In a hearing aid, it describes the amount of amplification provided for low, middle, and high pitches.

\medterm{Frequency, Urinary} The need to urinate more often than usual, sometimes passing only small amounts. Infection, irritation, pregnancy, diuretics, diabetes, overactive bladder, or prostate disease may cause it.

\medterm{Frequent Bowel Movements} Passing stool more often than is usual for a person. Causes include diet, infection, medicines, inflammation, poor absorption, or functional bowel disorders; blood, fever, dehydration, or weight loss needs assessment.

\medterm{Frequent Or Urgent Urination} Urinating often or having sudden urges that are difficult to postpone. Infection, bladder irritation, diabetes, pregnancy, medicines, obstruction, or overactive bladder can cause these symptoms.

\medterm{Frequent Urination} Needing to urinate more often than usual. It may result from drinking more, infection, diabetes, pregnancy, medicines, prostate enlargement, overactive bladder, or another urinary condition.

\medterm{Fresh Specimen} A recently collected biological sample tested promptly or preserved correctly before changes affect its cells, organisms, chemistry, or diagnostic accuracy.




\medterm{Frey’S Syndrome} Sweating, warmth, or flushing over the cheek or temple while eating, caused by nerve regrowth after parotid-gland surgery or injury. Symptoms may be treated with topical medicines, injections, or other specialist care.

\medterm{Friable} Describes tissue that is unusually soft, fragile, and easily torn. Friable tissue may bleed readily when touched and can occur with inflammation, infection, cancer, or poor healing.

\medterm{Friction} Rubbing between surfaces. Skin friction can cause irritation, blisters, or sores; reducing moisture, pressure, and rubbing helps prevent injury.

\medterm{Friedman Curve} An older graph describing cervical opening over time during labor. Modern labor care recognizes wider normal variation and does not use this historical curve alone to diagnose slow labor or decide on intervention.

\medterm{Friedreich Ataxia} An inherited disorder that progressively damages nerves and coordination. It can cause unsteady walking, weakness, reduced sensation, spinal curvature, heart muscle disease, and sometimes diabetes.

\medterm{Friedreich'S Ataxia} An inherited condition that gradually affects balance, coordination, sensation, and strength. Some people also develop heart muscle disease, spinal curvature, or diabetes and need regular neurologic and cardiac care.
\medterm{Friedreichs Ataxia Syndrome} An inherited disorder that gradually damages nerves and coordination. It can cause unsteady walking, weakness, reduced sensation, spinal curvature, heart muscle disease, and diabetes and usually requires long-term specialist care.

\medterm{Frigidity} An outdated and stigmatizing term for sexual difficulties. Modern assessment describes the specific problem, such as reduced desire, arousal, orgasm, or pain, and considers physical, emotional, relationship, and medicine-related causes.
\medterm{Frontal} Relating to the forehead, frontal skull bone, or frontal lobe of the brain. The frontal lobe helps control voluntary movement, planning, attention, speech, behavior, and decision-making.

\medterm{Frontal Bone} The skull bone forming the forehead, roof of the eye sockets, and part of the nasal region. It contains the frontal sinuses and protects the front of the brain.

\medterm{Frontal Bossing} An unusually prominent forehead caused by forward bulging of the frontal skull bones. It may be familial or associated with bone disease, excess pressure in the skull, or skeletal disorders.

\medterm{Frontal Lobe} The front portion of the brain’s outer layer. It helps control movement, planning, attention, judgment, behavior, personality, and speech production in the language-dominant hemisphere.

\medterm{Frontal Lobe Epilepsy} Epilepsy caused by seizures beginning in the frontal lobe. Episodes may involve sudden movements, unusual behavior, impaired awareness, or events during sleep and can be difficult to distinguish from sleep disorders.

\synonyms
Frontal lobe seizures

\medterm{Frontal Lobe Seizures} Seizures beginning in the front part of the brain, often causing brief unusual movements, behaviors, sensations, speech changes, or impaired awareness. Symptoms depend on the exact brain area involved.

\synonyms
Epilepsy, Frontal Lobe

\medterm{Frontal Sinus} An air-filled cavity within the frontal skull bone above the eye. It connects with the nasal passages and can become inflamed or blocked during sinusitis.

\medterm{Frontal Sinusitis} Inflammation of the air-filled cavities above the eyes, usually from a viral infection or blocked drainage. It can cause forehead pain, congestion, and discharge; severe headache, eye swelling, fever, confusion, or neurologic symptoms require urgent care.

\medterm{Frontotemporal Dementia} A group of progressive brain disorders affecting the frontal and temporal lobes. Early changes often involve personality, behavior, judgment, language, or movement rather than memory alone.

\synonyms
Dementia, Frontotemporal

\medterm{Frontotemporal Lobar Degeneration} Progressive loss of nerve cells in the frontal and temporal brain regions. It can cause early changes in behavior, personality, language, judgment, or movement and may occur at a relatively young age.

\medterm{Frostbite} A freezing injury to the skin and deeper tissues that can cause numbness, pale or waxy skin, swelling, blisters, and tissue death.

\medterm{Frotteuristic Disorder} A mental disorder involving persistent sexual arousal from touching or rubbing against a person without consent, causing distress or impairment or involving acting on the urges. Nonconsensual touching is sexual assault.

\medterm{Frozen Section} A rapid tissue examination performed during surgery. A specimen is frozen, cut into thin slices, and examined to provide immediate information about a lesion or surgical margin, although permanent sections are often more accurate.

\medterm{Frozen Shoulder} A painful condition in which the shoulder capsule tightens, progressively limiting movement. It often improves slowly over months or years; diabetes, thyroid disease, and prolonged immobility increase risk.

\synonyms
Adhesive Capsulitis Of Shoulder (Frozen Shoulder)


\medterm{Fructosamine} A blood measurement of sugar attached to blood proteins, reflecting average blood glucose over roughly the previous two to three weeks. It can help when the longer-term hemoglobin A1c test is unreliable.

\medterm{Fructose} A natural sugar found in fruit, honey, some vegetables, and many sweetened foods. It is processed mainly by the liver; poor absorption can cause gas, bloating, cramps, or diarrhea after eating it.

\synonyms
Fruit sugar

\medterm{Fructose Intolerance} Difficulty digesting or processing fructose, depending on the condition. Poor intestinal absorption causes gas, bloating, and diarrhea, while inherited inability to process fructose can cause low blood sugar and liver injury.

\medterm{Fructose Malabsorption} Incomplete absorption of fructose in the small intestine. Unabsorbed sugar is fermented by bacteria and can cause gas, bloating, cramps, or diarrhea; reducing personal dietary triggers may help.

\medterm{Fructose Metabolism} The body’s processing of fructose, mainly in the liver, into substances used for energy or stored. Rare enzyme defects can cause harmless fructose in urine or serious low blood sugar and liver injury after intake.


\medterm{Fruity Breath} A sweet, acetone-like breath odor caused by ketones. In someone with diabetes, especially with high glucose, vomiting, abdominal pain, deep breathing, or confusion, it may signal life-threatening ketoacidosis.


\medterm{Frustration} An emotional response to blocked goals, unmet needs, or difficult circumstances. It may cause irritability, anger, distress, withdrawal, or behavior changes and becomes clinically important when persistent or impairing.

\medterm{Fryns Syndrome} A rare inherited condition present from birth, often involving a hole in the diaphragm, characteristic facial features, and abnormal fingers or toes. Brain, heart, kidney, or genital abnormalities may also occur.

\medterm{FSH} An abbreviation for follicle-stimulating hormone, which supports ovarian follicle development and sperm production. Blood levels help assess infertility, puberty, menopause, and disorders of the pituitary gland, ovaries, or testes.

\synonyms
Follicle-stimulating hormone

\medterm{Ft.} An abbreviation for foot or feet, a length of 12 inches or about 30.48 centimeters. In medical records, the surrounding context should distinguish it from other abbreviations.

\medterm{FTA-ABS Blood Test} A blood test detecting antibodies against Treponema pallidum, the bacterium that causes syphilis. It supports diagnosis but often stays positive after treatment and cannot by itself show whether infection is active.

\medterm{FTA-ABS Test} A syphilis blood test that detects antibodies against Treponema pallidum. Because antibodies may remain for life after infection, clinicians interpret the result with other tests, symptoms, treatment history, and exposure risk.

\medterm{Fuchs Dystrophy} A progressive eye disorder in which cells lining the inner cornea gradually fail to remove fluid. It can cause blurred morning vision, glare, corneal swelling, pain, and worsening sight.

\medterm{Fuchs Endothelial Corneal Dystrophy} A usually bilateral disorder in which the cornea’s inner pumping cells are lost and tiny deposits develop. Fluid then accumulates, causing morning blur, glare, swelling, pain, and reduced vision.

\synonyms
Fuchs dystrophy

\medterm{Fuchs' Dystrophy} A progressive disorder of the cornea’s inner cells that causes fluid buildup, swelling, blurred vision, glare, and sometimes painful blisters on the corneal surface.
\medterm{Fucose} A sugar found in some cell-surface proteins and fats. It helps cells recognize and attach to one another and participates in immune signaling and other biological processes.

\medterm{Fucosidosis} A rare inherited disorder in which cells cannot break down certain sugar-containing substances. Their buildup progressively affects development, movement, bones, growth, sweating, and organ function.

\medterm{Fucosyltransferase} An enzyme that attaches the sugar fucose to proteins or fats. This changes how cells recognize one another, develop, interact with microbes, and display some blood-group markers.

\medterm{Fuel Oil Poisoning} Illness after swallowing, inhaling, or having fuel oil on the skin. Inhalation can damage the lungs and cause coughing or breathing difficulty; vomiting should not be induced, and urgent poison-control advice is needed.

\medterm{Fugue} A rare dissociative state involving unexpected travel or wandering with loss of memory about personal identity or the past. It needs assessment for trauma, mental illness, substances, neurologic disease, and safety risks.

\medterm{Fukuyama Congenital Muscular Dystrophy} A rare inherited muscle disorder beginning in infancy, causing low muscle tone, weakness, contractures, and delayed development. Abnormalities of the brain and eyes, seizures, and heart disease may also occur.

\synonyms
FCMD; Fukuyama Type Congenital Muscular Dystrophy

\medterm{Fulguration} Destruction or sealing of tissue using electrical energy during a procedure. It may control bleeding or treat selected abnormal tissue, but nearby tissue can be burned or injured.

\medterm{Full Agonist} A medicine or chemical that binds to a cell receptor and produces the strongest response that receptor can generate. The actual effect still depends on dose, receptor number, tissue, and individual response.

\medterm{Full Blood Count} See \textit{Complete Blood Count}.

\synonyms
FBC

\medterm{Full Liquid Diet} A diet containing liquids and foods that become liquid at room temperature, such as milk, strained soups, yogurt, and custard. It may be used temporarily, but prolonged use can lack fiber and some nutrients.

\medterm{Full Sibling} A person who shares both biological parents with another person. This relationship is relevant to family medical history, genetic-risk assessment, and some tissue or organ-donor matching.

\medterm{Full-Thickness Skin Graft} A graft containing the outer skin layer and the entire dermis, moved to cover a wound or reconstruct a defect. It usually gives durable coverage but requires a suitable donor site and blood supply.

\medterm{Fulminant} Beginning suddenly and becoming severe very quickly. The term describes a dangerous pattern of illness, such as rapidly progressive liver failure, infection, or inflammation requiring urgent treatment.

\medterm{Fulminant Hepatitis} Rapid, severe liver injury causing loss of liver function, abnormal bleeding, and sometimes confusion or coma from toxin buildup. It is a medical emergency that may require intensive care and liver transplantation.

\medterm{Fulminant Myocarditis} Sudden, severe inflammation of the heart muscle causing poor pumping, dangerous rhythms, or shock. It requires intensive care and treatment of the cause and complications; temporary mechanical heart support may be needed.

\medterm{Fulvestrant} An injectable hormone treatment that blocks and helps destroy estrogen receptors in selected hormone-sensitive breast cancers. It can cause injection-site pain, hot flushes, joint symptoms, fatigue, and liver-test changes.
\medterm{Fumigation} Treating an enclosed area with gas, vapor, or smoke to kill pests or microorganisms. People and animals must leave beforehand and return only after the area is declared safe and adequately ventilated.

\medterm{Functional Asplenia} Poor or absent spleen function even though the spleen is present. It increases the risk of severe infections, especially from certain bacteria, so vaccination, prompt fever assessment, and sometimes preventive antibiotics are important.

\medterm{Functional Assessment} Evaluation of how well a person moves, communicates, thinks, and performs daily activities. It identifies support needs, tracks change, guides rehabilitation, and helps plan safe care.

\medterm{Functional Capacity} The level of physical, mental, communication, and self-care activity a person can safely perform under particular conditions. It depends on health, environment, assistance, and personal goals.

\medterm{Functional Disorder} A condition in which symptoms or impaired body function arise from altered processing or regulation without a structural abnormality sufficient to explain them. The symptoms are real and can cause substantial disability.

\medterm{Functional Dyspepsia} Ongoing upper-abdominal discomfort, burning, early fullness, or unpleasant fullness after meals without an ulcer or other structural cause explaining symptoms. It may involve altered stomach movement or gut sensitivity.

\synonyms
Nonulcer Dyspepsia
Nonulcer Stomach Pain

\medterm{Functional Electrical Stimulation} Electrical impulses applied to nerves or muscles to produce a useful movement, such as lifting the foot during walking. It may assist rehabilitation after stroke, spinal injury, or other nerve damage.

\medterm{Functional Gastrointestinal Disorders} Digestive symptoms caused by altered gut movement or sensitivity without an obvious structural disease. Symptoms may include pain, bloating, constipation, diarrhea, nausea, or early fullness.

\synonyms
Disorders of gut-brain interaction

\medterm{Functional Hypothalamic Amenorrhea} Missing menstrual periods because stress, inadequate energy intake, excessive exercise, or disordered eating suppresses reproductive hormones. It can lower estrogen, impair fertility, and weaken bones, but often improves when the underlying imbalance is corrected.

\synonyms
Hypothalamic Amenorrhea; FHA

\medterm{Functional Incontinence} Urine leakage because a person cannot reach, recognize, or use the toilet in time, even when bladder storage may be adequate. Mobility, memory, vision, clothing, communication, or environmental barriers may contribute.

\medterm{Functional Independence Measure} A structured tool that rates how much help a person needs with self-care, toileting, transfers, walking, communication, and thinking. It is used mainly in rehabilitation to track function and progress.

\medterm{Functional MRI} A magnetic-resonance scan that estimates brain activity from changes in blood flow and oxygen use. It can map language or movement areas before selected brain procedures and does not use X-ray radiation.

\medterm{Functional Murmur} An extra heart sound caused by increased blood flow through a structurally normal heart. It may occur with fever, anemia, pregnancy, or exercise and is assessed to exclude valve or heart disease.

\medterm{Functional Neurologic Disorder} A disorder in which altered nervous-system function causes genuine symptoms such as weakness, abnormal movements, seizures, or sensory changes without the structural damage expected for them.
\synonyms
FND; functional neurological disorder; functional neurological symptom disorder; conversion disorder

\medterm{Functional Neurologic Disorder/Conversion Disorder} Alternate index label; see \textit{Functional Neurologic Disorder}.

\medterm{Functional Neurological Disorder} Alternate index label; see \textit{Functional Neurologic Disorder}.

\medterm{Functional Neurological Symptom Disorder} Alternate index label; see \textit{Functional Neurologic Disorder}.

\medterm{Functional Neurological Symptom Disorder (Conversion Disorder)} Alternate index label; see \textit{Functional Neurologic Disorder}.

\medterm{Functional Residual Capacity} The amount of air left in the lungs after a normal relaxed breath out. It helps keep air sacs open and may be reduced by obesity or restrictive lung disease or increased by air trapping.

\medterm{Functional Screening} A brief check of a person’s ability to move safely, communicate, perform daily activities, and manage basic self-care. Abnormal results may lead to a fuller medical, occupational, or rehabilitation assessment.

\medterm{Functional Status} A person’s ability to manage daily activities such as bathing, dressing, eating, toileting, walking, shopping, cooking, medicines, finances, and transportation. Decline may signal illness or a need for support.

\medterm{Functional Urinary Incontinence} Urine leakage that happens because a person cannot reach, recognize, or use the toilet in time, even when the bladder may work normally. Problems with walking, memory, vision, clothing, communication, or the environment can contribute; treatment addresses these barriers and other causes of leakage.

\medterm{Fundal Height} The distance from the pubic bone to the top of the uterus, measured during pregnancy. It roughly follows weeks of pregnancy after about 20 weeks; differences may reflect dating error, growth problems, twins, or fluid changes.

\medterm{Fundoplication} Surgery that wraps the upper stomach around the lower esophagus to strengthen the valve and reduce acid reflux. It may also repair a hiatal hernia; swallowing difficulty, gas, or recurrent reflux can occur afterward.

\medterm{Fundoscopy} Examination of the back of the eye using an ophthalmoscope or camera. It views the retina, optic nerve, macula, and blood vessels to assess conditions such as diabetes, high blood pressure, or retinal disease.

\medterm{Fundus} The rounded or farthest part of a hollow organ. Examples include the upper part of the stomach, the top of the uterus, and the inside back surface of the eye.

\medterm{Fundus Oculi} The inside back surface of the eye visible through the pupil. It includes the retina, optic nerve, macula, and retinal blood vessels.

\medterm{Fundus Photography} Photographing the back of the eye to record the retina, optic nerve, macula, and blood vessels. Images help detect and monitor retinal disease and allow comparison over time.

\medterm{Fundus Uteri} The rounded upper part of the uterus above the openings of the fallopian tubes. Its position and size can be assessed during pelvic examination and pregnancy.

\medterm{Fundus, Retinal} The inside back surface of the eye viewed during fundoscopy, including the retina, optic nerve, macula, and retinal blood vessels.

\medterm{Fungal} Relating to fungi or caused by a fungus. Fungal infections may affect the skin, nails, mouth, lungs, bloodstream, or other organs.

\medterm{Fungal Arthritis} A fungal infection inside a joint, usually after surgery, injury, injection, bloodstream spread, or with weakened immunity. It often develops slowly with persistent pain, swelling, and reduced movement and may need prolonged treatment and surgery.

\medterm{Fungal Endophthalmitis} A fungal infection inside the eye, often after bloodstream infection, surgery, injury, or weakened immunity. It can rapidly damage vision and requires urgent specialist treatment with antifungal medicine and sometimes surgery.

\medterm{Fungal Infection} An infection caused by a fungus. It may affect the skin, nails, or mouth, or spread to the lungs, bloodstream, or other organs, especially when the immune system is weakened.

\medterm{Fungal Infection (Tinea/Dermatophytosis)} A superficial infection of skin, hair, or nails caused by dermatophyte fungi. Ringworm, athlete’s foot, and jock itch commonly cause itchy, scaly, or ring-shaped rashes in warm, moist areas.

\synonyms
Tinea/Dermatophytosis

\medterm{Fungal Infection, Nail} A fungal infection of a fingernail or toenail causing thickening, discoloration, brittleness, or separation from the nail bed. Diagnosis may require testing because injury and skin disease can look similar.

\synonyms
Onychomycosis; nail fungus

\medterm{Fungal Nail Infection} Infection of a fingernail or toenail by a fungus, causing thickening, discoloration, brittleness, debris beneath the nail, or separation from its bed. Treatment depends on the organism, severity, and medical history.

\medterm{Fungal Nails} Nails infected by fungi may become thick, yellow or brown, brittle, crumbly, or separated from the nail bed. Testing may be needed before treatment, which can be topical or oral.

\synonyms
Onychomycosis (Fungal Nails)

\medterm{Fungal Pathogenesis} The ways fungi cause disease, including attaching to cells, growing into tissue, avoiding immune defenses, and releasing damaging substances. The process varies with the fungus, body site, and person’s immunity.

\medterm{Fungal Resistance Mechanisms} Changes that allow fungi to survive antifungal medicines, such as altering a drug target, pumping medicine out of the cell, or changing cell-wall production. Resistance can make infections harder to treat.

\medterm{Fungal Skin Infections} Infections of the skin caused by dermatophytes or yeasts, including ringworm, athlete’s foot, and candidiasis. They often cause red, itchy, scaly, or ring-shaped areas and thrive in warm, moist conditions.

\medterm{Fungating} Describing a mass or wound that grows outward with an irregular, ulcerated, dead, or foul-smelling surface. It may occur with advanced cancer or severe chronic wounds and can bleed, leak, and become infected.

\medterm{Fungemia} The presence of living fungi in the bloodstream, most often Candida in hospitalized patients. It can spread to organs and requires urgent cultures, source evaluation, and appropriate antifungal treatment.

\medterm{Fungi} Organisms including yeasts, molds, and mushrooms. Some live harmlessly in the environment or body, while others cause infections of the skin, mouth, lungs, bloodstream, or other organs.

\medterm{Fungicide} A chemical used to kill or inhibit fungi. Some fungicides are used on crops or surfaces, while selected medicines act against fungi in people; toxicity depends on the substance and exposure.

\medterm{Fungoid} Resembling a fungus or mushroom in appearance. In medicine it may describe an outward-growing, irregular lesion or mass, but appearance alone does not establish its cause.

\medterm{Fungus} A type of organism including yeasts and molds, with distinctive cell walls and membranes. Fungi reproduce in various ways and can cause superficial or invasive infection, especially when immunity is weakened.

\medterm{Funiculus} A cord-like bundle of nerve fibers or other structures. In the spinal cord, a funiculus contains nerve pathways; in the spermatic cord, it is the bundle carrying the vas deferens, blood vessels, nerves, and lymph vessels.

\medterm{Funisitis} Inflammation of the umbilical cord, usually associated with infection or inflammation of the membranes around the fetus. It may increase newborn infection risk and requires assessment of the placenta, mother, and baby.

\medterm{Funnel Chest} A chest-wall shape present from birth in which the breastbone and nearby ribs curve inward. It is often mild, but severe deformity may affect exercise, breathing, body image, or heart function.

\synonyms
Pectus excavatum

\medterm{Funnel-Web Spider Bite} A potentially life-threatening bite from a funnel-web spider. Venom may cause pain, sweating, tingling, muscle spasms, breathing difficulty, or collapse; keep the person still, apply pressure immobilization, and seek emergency care.

\medterm{Furan} A chemical compound used in manufacturing and found in small amounts in some heated foods. High or repeated exposure may harm the liver and other organs, so occupational exposure should be controlled.

\medterm{Furniture Polish Poisoning} Illness after swallowing or inhaling furniture-polish chemicals. Some products can damage the lungs if aspirated and cause drowsiness or chemical injury; do not induce vomiting and contact poison-control services urgently.

\medterm{Furocoumarin} A plant chemical found in some citrus and other plants that can make skin unusually sensitive to ultraviolet light. Contact followed by sunlight may cause redness, blistering, or dark pigmentation.

\medterm{Fursultiamin} A fat-soluble form of thiamine, or vitamin B1, used in some countries as a vitamin preparation. Its clinical value and safety depend on the indication, dose, and other medicines.

\medterm{Furuncle} A painful, pus-filled infection of a hair follicle, usually caused by Staphylococcus aureus. It forms a red lump and may drain; larger, recurrent, facial, or fever-associated lesions need medical assessment.

\medterm{Furuncle And Carbuncle} A furuncle is a deep infected hair follicle, or boil. A carbuncle is a connected group of boils forming a larger, deeper infection; both may require drainage and sometimes antibiotics.

\medterm{Furuncular Myiasis} A fly-larva infestation forming a boil-like skin lump with a central opening through which the larva breathes. It may cause pain, movement sensations, or drainage and requires safe removal.

\medterm{Furunculosis} Repeated or multiple boils caused by infection of hair follicles, often with Staphylococcus aureus. Risk factors include close contact, skin disease, diabetes, and weakened immunity; treatment may require drainage, antibiotics, and prevention measures.

\medterm{Fusariosis} An opportunistic infection caused by Fusarium fungi, mainly affecting people with severely weakened immunity or prolonged low neutrophil counts. It can cause skin lesions, bloodstream infection, pneumonia, and disease in multiple organs.

\medterm{Fusarium} A group of environmental molds that can cause nail, eye, skin, sinus, lung, or widespread infections. Severe disease occurs mainly in people with weakened immunity, especially prolonged neutropenia.

\medterm{Fusarium Infection} An infection caused by Fusarium molds, which may affect the nails, eyes, skin, sinuses, lungs, or bloodstream. Invasive disease mainly affects people with severe immune suppression and requires urgent specialist antifungal care.
\medterm{Fusarium Solani} An environmental mold that can cause eye, nail, skin, or invasive infections, especially in people with weakened immunity. Its clinical importance depends on the specimen, body site, symptoms, and laboratory confirmation.

\medterm{Fused Placenta} A placenta in which areas that normally form separate lobes are joined. It is often an anatomical variation, but careful examination after delivery helps identify any retained placental tissue or abnormal blood vessels.

\medterm{Fused Teeth} Alternate index label for dental fusion.

\medterm{Fusidic Acid} An antibiotic that blocks bacterial protein production and is used for selected staphylococcal infections. Resistance can develop, so it is often combined with another antibiotic and used according to local guidance.

\medterm{Fusiform} Shaped like a spindle, wider in the middle and narrowing toward both ends. The term may describe a muscle, cell, swelling, or blood-vessel enlargement such as an aneurysm.

\medterm{Fusion} Joining of two structures that are normally separate. It may describe bones, joints, teeth, organs, genes, or tissues and can occur naturally, during development, after injury, or through surgery.

\synonyms
Dental fusion
Tooth fusion

\medterm{Fusion Gene} A gene formed when parts of two separate genes become joined. It may produce an abnormal protein that drives some cancers and can sometimes be identified by genetic testing.

\medterm{Fusion Of The Ear Bones} Abnormal joining or stiffening of the tiny middle-ear bones, which can reduce sound transmission and cause conductive hearing loss. Evaluation identifies the cause and whether surgery or hearing support may help.
\medterm{Fusion Protein} A protein made when parts of two different proteins become joined, usually because their genes have fused. Some fusion proteins drive cancer; others are deliberately engineered for research or treatment.

\medterm{Fusobacterium} A group of bacteria that grow without oxygen and commonly live in the mouth and intestine. Some species can cause dental, throat, abdominal, wound, bloodstream, or other invasive infections.

\medterm{Fusobacterium Infection} Infection caused by Fusobacterium bacteria, which may spread from the mouth or intestine to the throat, lungs, abdomen, skin, joints, or bloodstream. Symptoms and treatment depend on the site and severity.

\medterm{Fusobacterium Necrophorum} A bacterium that can cause severe throat infection and Lemierre syndrome, in which infection spreads to neck veins and may produce infected clots, lung complications, and bloodstream infection.

\medterm{Fusobacterium Nucleatum} A common mouth bacterium associated with gum disease. It can occasionally cause bloodstream or other invasive infection and has also been studied for possible links with intestinal and pregnancy-related disease.
\medterm{Fusospirochetal Gingivitis} A painful gum infection involving fusobacteria and spirochetes, causing bleeding, ulcers, swelling, and foul breath. It is associated with poor oral hygiene, stress, smoking, or illness and needs dental treatment.
\medterm{Fibular Artery} Also called the peroneal artery, it is a branch of the leg’s main artery that supplies the outer and back lower leg and contributes blood flow around the ankle.

\medterm{Fahr's Syndrome} A rare disorder involving calcium deposits in deep brain regions, sometimes causing movement problems, seizures, thinking changes, or psychiatric symptoms. Doctors look for inherited causes and treatable problems with calcium or parathyroid hormones.

\medterm{Fallen Bladder} A common term for bladder prolapse, in which weakened pelvic-support tissues allow the bladder to bulge into the front wall of the vagina. It may cause pressure, leakage, incomplete emptying, or no symptoms.

\medterm{Fatty Liver of Pregnancy (AFLP)} A rare, serious liver disorder usually occurring late in pregnancy. It can cause nausea, abdominal pain, jaundice, low blood sugar, abnormal bleeding, kidney injury, or liver failure and requires urgent delivery planning and specialist care.

\medterm{Fecal Transplant} A procedure placing carefully screened donor stool into the intestine to restore helpful gut microbes. It is mainly used for selected recurrent C. difficile infection and has small but important infection and procedure-related risks.

\medterm{Femoral Anteversion} An inward twist of the thigh bone that turns the knees or feet inward, often noticed in childhood. It commonly improves with growth; persistent pain, severe asymmetry, or walking problems need orthopedic assessment.

\medterm{Fertility Treatment} Medical or procedural care intended to help achieve pregnancy. Options include treating underlying disease, ovulation medicines, surgery, insemination, or in-vitro fertilization, chosen according to cause, age, health, and reproductive goals.

\medterm{Fetal Alcohol Spectrum Disorders (FASD)} A group of lifelong conditions caused by alcohol exposure before birth. Effects may include poor growth, learning or attention problems, behavior changes, coordination difficulties, or characteristic facial features; no amount of alcohol is known to be safe during pregnancy.

\medterm{Fetor Hepaticus} A musty or sweet breath odor that can occur with severe liver dysfunction when certain substances are not adequately cleared. It is a warning sign requiring evaluation for advanced liver disease and complications.

\medterm{Fish Odour Syndrome} A rare metabolic disorder in which trimethylamine builds up and causes a fish-like smell in sweat, breath, or urine. Symptoms may vary with diet; testing, dietary guidance, and emotional support can help.

\medterm{Focal Segmental Glomerular Sclerosis} A kidney disorder in which scars affect some filters or parts of them. It can cause protein in urine, swelling, high blood pressure, and declining kidney function and may result from genetic, immune, or other conditions.

\medterm{FODMAP} An abbreviation for fermentable short-chain carbohydrates that may be poorly absorbed and cause gas, bloating, pain, or diarrhea. A supervised reduction followed by gradual reintroduction may help some people with irritable bowel syndrome.

\medterm{Folliculitis Decalvans} A rare inflammatory scalp disorder causing repeated pustules, crusting, hair loss, and permanent scarring. Early dermatologic treatment aims to control inflammation and infection and limit further follicle destruction.

\medterm{Food Addiction} A nonformal term for intense cravings or loss of control around eating. It is not a universally accepted diagnosis; assessment should consider binge-eating disorder, mood, stress, medicines, sleep, and access to food.

\medterm{Food Contamination} The presence of harmful germs, toxins, chemicals, allergens, or foreign materials in food. Contamination can occur during production, storage, preparation, or serving and may cause infection, poisoning, allergic reactions, or injury.

\medterm{Forensics} The use of scientific methods to answer legal questions, including identification, toxicology, injury analysis, and interpretation of biological or physical evidence. Reliable findings require validated methods, documentation, and expert interpretation.

\medterm{Fowler's Syndrome} A rare cause of persistent difficulty emptying the bladder, usually in young women, related to abnormal relaxation of the urethral sphincter. Evaluation excludes blockage and nerve disease; treatment may require drainage or neuromodulation.

\medterm{Fox-Fordyce Disease} A rare chronic skin disorder caused by blocked and inflamed sweat glands. It produces intensely itchy small bumps, commonly in the armpits, pubic region, or around the nipples.

\medterm{Frontal Lobe Disorder} A condition affecting brain areas involved in planning, judgment, movement, attention, speech, and behavior. Causes include stroke, injury, tumor, dementia, infection, or other neurologic disease.

\medterm{Frontalis Suspension} An operation connecting the upper eyelid to the forehead muscle with a sling to lift a drooping eyelid. It is used when the eyelid-lifting muscle is weak, especially in congenital or childhood ptosis.
